
%\newpage
\section{名词解释-part1.tex should be deleted}

\subsection{忍耐}
\subsubsection{路加福音8:15/21:14-19}
\begin{itemize}[noitemsep]
\item 8:15 那落在好土里的，就是人听了道，持守在诚实善良的心里，并且\textcolor{red}{忍耐}着结实。
\item 21:14-19 所以，你们当立定心意，不要预先思想怎样分诉；因为我必赐你们口才、智慧，是你们一切敌人所敌不住、驳不倒的。连你们的父母、弟兄、亲族、朋友也要把你们交官，你们也有被他们害死的。你们要为我的名被众人恨恶，然而，你们连一根头发也必不损坏。你们常存\textcolor{red}{忍耐}，就必保全灵魂(必得生命)。
\end{itemize}


\subsubsection{罗马书2:7/5:3-4/8:25/15:4-5} 
\begin{itemize}[noitemsep]
\item 2:7 他必照各人的行为报应各人。凡\textcolor{red}{恒心(耐心)}行善寻求荣耀、尊贵和不能朽坏之福的,就以永生报应他们
\item 5:3-4 我们既因信称义，就借着我们的主耶稣基督得与神相和。 2 我们又借着他，因信得进入现在所站的这恩典中，并且欢欢喜喜盼望神的荣耀。 3 不但如此，就是在患难中也是欢欢喜喜的。因为知道患难生\textcolor{red}{忍耐}， 4 \textcolor{red}{忍耐}生老练，老练生盼望， 5 盼望不至于羞耻，因为所赐给我们的圣灵将神的爱浇灌在我们心里。
\item 8:25 我们得救是在乎盼望。只是所见的盼望不是盼望。谁还盼望他所见的呢？25 但我们若盼望那所不见的，就必\textcolor{red}{忍耐}等候。
\item 15:4-5 从前所写的圣经都是为教训我们写的，叫我们因圣经所生的\textcolor{red}{忍耐}和安慰，可以得着盼望。 5 但愿赐\textcolor{red}{忍耐}、安慰的神叫你们彼此同心，效法基督耶稣， 6 一心一口荣耀神我们主耶稣基督的父！
\end{itemize}


\subsubsection{哥林多后书 1:6/6:4/12:12} 
\begin{itemize}[noitemsep]
\item 1:6 我们在一切患难中，他就安慰我们，叫我们能用神所赐的安慰去安慰那遭各样患难的人。 5 我们既多受基督的苦楚，就靠基督多得安慰。 6 我们受患难呢，是为叫你们得安慰、得拯救；我们得安慰呢，也是为叫你们得安慰。这安慰能叫你们\textcolor{red}{忍受}我们所受的那样苦楚。 7 我们为你们所存的盼望是确定的，因为知道你们既是同受苦楚，也必同得安慰。
\item 6:4 我们凡事都不叫人有妨碍，免得这职分被人毁谤；反倒在各样的事上表明自己是神的用人，就如在许多的\textcolor{red}{忍耐}，患难，穷乏，困苦，鞭打，监禁，扰乱，勤劳，警醒，不食，廉洁，知识，恒忍，恩慈，圣灵的感化，无伪的爱心，真实的道理，神的大能；仁义的兵器在左在右，荣耀羞辱，恶名美名；似乎是诱惑人的，却是诚实的；似乎不为人所知，却是人所共知的；似乎要死，却是活着的；似乎受责罚，却是不致丧命的；似乎忧愁，却是常常快乐的；似乎贫穷，却是叫许多人富足的；似乎一无所有，却是样样都有的。
\item 12:12 我成了愚妄人，是被你们强逼的。我本该被你们称许才是。我虽算不了什么，却没有一件事在那些最大的使徒以下。 12 我在你们中间用百般的\textcolor{red}{忍耐}，借着神迹、奇事、异能，显出使徒的凭据来。 13 除了我不累着你们这一件事，你们还有什么事不及别的教会呢？这不公之处，求你们饶恕我吧！
\end{itemize}



\subsubsection{帖前 1:3/帖后 1:4/3:5} 
\begin{itemize}[noitemsep]
\item 帖前 1:3 我们为你们众人常常感谢神，祷告的时候提到你们， 3 在神我们的父面前不住地记念你们因信心所做的工夫，因爱心所受的劳苦，因盼望我们主耶稣基督所存的\textcolor{red}{忍耐}。
\item 帖后 1:4 弟兄们，我们该为你们常常感谢神，这本是合宜的，因你们的信心格外增长，并且你们众人彼此相爱的心也都充足。 4 甚至我们在神的各教会里为你们夸口，都因你们在所受的一切逼迫患难中，仍旧存\textcolor{red}{忍耐}和信心。
\item 3:5 主是信实的，要坚固你们，保护你们脱离那恶者(脫離凶惡)。 4 我们靠主深信，你们现在是遵行我们所吩咐的，后来也必要遵行。 5 愿主引导你们的心，叫你们爱神，并学基督的\textcolor{red}{忍耐}！
\end{itemize}

 

\subsubsection{提前 6:11/提后 3:10/提多书 2:2} 
\begin{itemize}[noitemsep]
\item 提前 6:11 那些想要发财的人，就陷在迷惑，落在网罗和许多无知有害的私欲里，叫人沉在败坏和灭亡中。 10 贪财是万恶之根。有人贪恋钱财，就被引诱离了真道，用许多愁苦把自己刺透了。但你这属神的人要逃避这些事，追求公义、敬虔、信心、爱心、\textcolor{red}{忍耐}、温柔。 
\item 提后 3:10 从前雅尼和佯庇怎样敌挡摩西，这等人也怎样敌挡真道。他们的心地坏了，在真道上是可废弃的。 9 然而他们不能再这样敌挡，因为他们的愚昧必在众人面前显露出来，像那二人一样。 10 但你已经服从了我的教训、品行、志向、信心、宽容、爱心、\textcolor{red}{忍耐}， 11 以及我在安提阿、以哥念、路司得所遭遇的逼迫、苦难。我所忍受是何等的逼迫，但从这一切苦难中，主都把我救出来了。 12 不但如此，凡立志在基督耶稣里敬虔度日的，也都要受逼迫。 13 只是作恶的和迷惑人的必越久越恶，他欺哄人，也被人欺哄。 14 但你所学习的、所确信的，要存在心里，因为你知道是跟谁学的， 15 并且知道你是从小明白圣经，这圣经能使你因信基督耶稣有得救的智慧。
\item 提多 2:2 劝老年人要有节制，端庄，自守，在信心、爱心、\textcolor{red}{忍耐}上都要纯全无疵。
\end{itemize}




\subsubsection{希伯来书 10:36/12:1} 
\begin{itemize}[noitemsep]
\item 10:36 你们要追念往日蒙了光照以后，所忍受大争战的各样苦难。一面被毁谤，遭患难，成了戏景叫众人观看；一面陪伴那些受这样苦难的人。因为你们体恤了那些被捆锁的人，并且你们的家业被人抢去，也甘心忍受，知道自己有更美、长存的家业。 所以你们不可丢弃勇敢的心，存这样的心必得大赏赐。你们必须\textcolor{red}{忍耐},使你们行完了神的旨意,就可以得着所应许的。 37 “因为还有一点点时候，那要来的就来，并不迟延。

\item 12:1 我们既有这许多的见证人，如同云彩围着我们，就当放下各样的重担，脱去容易缠累我们的罪，存心\textcolor{red}{忍耐}，奔那摆在我们前头的路程， 2 仰望为我们信心创始成终的耶稣(仰望那将真道创始成终的耶稣。)。他因那摆在前面的喜乐，就轻看羞辱，忍受了十字架的苦难，便坐在神宝座的右边。 3 那忍受罪人这样顶撞的，你们要思想，免得疲倦灰心。 4 你们与罪恶相争，还没有抵挡到流血的地步。 5 你们又忘了那劝你们如同劝儿子的话说：“我儿，你不可轻看主的管教，被他责备的时候也不可灰心。 6 因为主所爱的，他必管教，又鞭打凡所收纳的儿子。” 7 你们所忍受的，是神管教你们，待你们如同待儿子。焉有儿子不被父亲管教的呢？ 8 管教原是众子所共受的，你们若不受管教，就是私子，不是儿子了。
\end{itemize}



\subsubsection{雅各书 1:3/5:10-11} 
\begin{itemize}[noitemsep]
\item 1:3 我的弟兄们，你们落在百般试炼中，都要以为大喜乐， 3 因为知道，你们的信心经过试验就生\textcolor{red}{忍耐}。 4 但\textcolor{red}{忍耐}也当成功，使你们成全、完备，毫无缺欠。

\item 5:10-11 弟兄们，你们不要彼此埋怨，免得受审判。看哪，审判的主站在门前了！ 10 弟兄们，你们要把那先前奉主名说话的众先知当做能受苦、能\textcolor{red}{忍耐}的榜样。 11 那先前\textcolor{red}{忍耐}的人，我们称他们是有福的。你们听见过约伯的\textcolor{red}{忍耐}，也知道主给他的结局，明显主是满心怜悯、大有慈悲。
\end{itemize}




\subsubsection{歌罗西书 1:11/彼得后书 1:6/啟示錄 3:10-11} 
\begin{itemize}[noitemsep]
\item 歌罗西书 1:11 我们自从听见的日子，也就为你们不住地祷告祈求，愿你们在一切属灵的智慧悟性上，满心知道神的旨意， 10 好叫你们行事为人对得起主，凡事蒙他喜悦，在一切善事上结果子，渐渐地多知道神； 11 照他荣耀的权能，得以在各样的力上加力，好叫你们凡事欢欢喜喜地\textcolor{red}{忍耐}宽容。
 
\item 彼得后书 1:6 你们要分外地殷勤：有了信心，又要加上德行；有了德行，又要加上知识； 6 有了知识，又要加上节制；有了节制，又要加上\textcolor{red}{忍耐}；有了\textcolor{red}{忍耐}，又要加上虔敬； 7 有了虔敬，又要加上爱弟兄的心；有了爱弟兄的心，又要加上爱众人的心。

\item 啟示錄 3:10-11 你既遵守我\textcolor{red}{忍耐}的道，我必在普天下人受試煉的時候，保守你免去你的試煉。 11我必快來！你要持守你所有的，免得人奪去你的冠冕。
\end{itemize}




\subsection{鑰匙}

\subsubsection{大衛家的鑰匙-以賽亞書 22:22; 啟示錄 3:7}
\begin{itemize} 
\itemsep-.25em
\item 我必將\underline{大衛家的鑰匙}放在他肩頭上，他開無人能關，他關無人能開。
\item 「你要寫信給非拉鐵非教會的使者說：『那聖潔、真實，拿著\underline{大衛的鑰匙}，開了就沒有人能關、關了就沒有人能開的說：我知道你的行為，你略有一點力量，也曾遵守我的道，沒有棄絕我的名，看哪，我在你面前給你一個敞開的門，是無人能關的。
\end{itemize} 

\subsubsection{天國的鑰匙-馬太福音 16:19}
我要把\underline{天國的鑰匙}給你，凡你在地上所捆綁的，在天上也要捆綁；凡你在地上所釋放的，在天上也要釋放。」

\subsubsection{知識的鑰匙-路加福音 11:52}
你們律法師有禍了！因為你們把\underline{知識的鑰匙}奪了去，自己不進去，正要進去的人你們也阻擋他們。


\subsubsection{死亡和陰間的鑰匙(無底坑)-啟示錄 1:18,9:1,20:1}
\begin{itemize} 
\itemsep-.25em
\item 又是那存活的；我曾死過，現在又活了，直活到永永遠遠，並且拿著\underline{死亡和陰間的鑰匙} 
\item 第五位天使吹號，我就看見一個星從天落到地上，有\underline{無底坑的鑰匙}賜給它。它開了無底坑，便有煙從坑裡往上冒，好像大火爐的煙，日頭和天空都因這煙昏暗了
\item 我又看見一位天使從天降下，手裡拿著\underline{無底坑的鑰匙}和一條大鏈子。 2他捉住那龍，就是古蛇，又叫魔鬼，也叫撒旦，把牠捆綁一千年
\end{itemize} 


\subsection{七靈,七個金燈臺,七星,七眼}
\subsubsection{出埃及記 25:37，37:23}
\begin{itemize} 
\itemsep-.25em
\item 要做燈臺的七個燈盞，祭司要點這燈，使燈光對照(出 25:37)。
\item 用精金做燈臺的七個燈盞，並燈臺的蠟剪和蠟花盤(出 37:23)。
%\item 我說：「我看見了一個純金的燈臺，頂上有燈盞，燈臺上有七盞燈，每盞有七個管子(撒迦利亞書 4:2)。
\item 我又看見寶座與四活物並長老之中有羔羊站立，像是被殺過的，有七角七眼，就是神的七靈，奉差遣往普天下去的(啟示錄 5:6)。
\end{itemize} 

\subsubsection{撒迦利亞書 3:8-10，4:6-10}
\begin{itemize} 
\itemsep-.25em
\item 「大祭司約書亞啊，你和坐在你面前的同伴都當聽（他們是做預兆的）：我必使我僕人大衛的苗裔發出。 9 看哪，我在約書亞面前所立的石頭，在一塊石頭上有七眼。萬軍之耶和華說：我要親自雕刻這石頭，並要在一日之間除掉這地的罪孽。 10 當那日，你們各人要請鄰舍坐在葡萄樹和無花果樹下。這是萬軍之耶和華說的。」
\item 他對我說：「這是耶和華指示所羅巴伯的。萬軍之耶和華說：不是倚靠勢力，不是倚靠才能，乃是倚靠我的靈方能成事。 7 大山哪，你算什麼呢？在所羅巴伯面前，你必成為平地。他必搬出一塊石頭，安在殿頂上，人且大聲歡呼說：『願恩惠，恩惠歸於這殿！』」 8 耶和華的話又臨到我說： 9 「所羅巴伯的手立了這殿的根基，他的手也必完成這工，你就知道萬軍之耶和華差遣我到你們這裡來了。 10 誰藐視這日的事為小呢？這七眼乃是耶和華的眼睛，遍察全地，見所羅巴伯手拿線砣就歡喜。」
\end{itemize} 

\subsubsection{啟示錄 1:4-6，12-20，2:1-3，5，3:1-2， 5:6}
\begin{itemize} 
\itemsep-.25em
\item 約翰寫信給亞細亞的七個教會。但願從那昔在、今在、以後永在的神和他寶座前的七靈， 5 並那誠實作見證的、從死裡首先復活、為世上君王元首的耶穌基督，有恩惠、平安歸於你們！他愛我們，用自己的血使我們脫離罪惡， 6 又使我們成為國民，做他父神的祭司。但願榮耀、權能歸給他，直到永永遠遠！阿們。
\item 我轉過身來，要看是誰發聲與我說話；既轉過來，就看見七個金燈臺， 13 燈臺中間有一位好像人子，身穿長衣，直垂到腳，胸間束著金帶。 14 他的頭與髮皆白，如白羊毛，如雪，眼目如同火焰， 15 腳好像在爐中鍛煉光明的銅，聲音如同眾水的聲音。 16 他右手拿著七星，從他口中出來一把兩刃的利劍，面貌如同烈日放光。17 我一看見，就仆倒在他腳前，像死了一樣。他用右手按著我說：「不要懼怕。我是首先的，我是末後的， 18 又是那存活的；我曾死過，現在又活了，直活到永永遠遠，並且拿著死亡和陰間的鑰匙。 19 所以你要把所看見的和現在的事，並將來必成的事，都寫出來。 20 論到你所看見在我右手中的七星和七個金燈臺的奧祕，那七星就是七個教會的使者，七燈臺就是七個教會。
\item 1 「你要寫信給以弗所教會的使者說：『那右手拿著七星、在七個金燈臺中間行走的說： 2 我知道你的行為、勞碌、忍耐，也知道你不能容忍惡人。你也曾試驗那自稱為使徒卻不是使徒的，看出他們是假的來。 3 你也能忍耐，曾為我的名勞苦，並不乏倦。5 所以，應當回想你是從哪裡墜落的，並要悔改，行起初所行的事。你若不悔改，我就臨到你那裡，把你的燈臺從原處挪去。
\item 「你要寫信給撒狄教會的使者說：『那有神的七靈和七星的說：我知道你的行為，按名你是活的，其實是死的。 2 你要警醒，堅固那剩下將要衰微的，因我見你的行為在我神面前，沒有一樣是完全的。
\item 我又看見寶座與四活物並長老之中有羔羊站立，像是被殺過的，有七角七眼，就是神的七靈，奉差遣往普天下去的。
\end{itemize} 

\subsection{白马}
\subsubsection{白马(啟 3:10, 6:1-2, 19:11-21/撒迦利亞書 1:7-11, 6:1-8)}
\begin{itemize}
\itemsep-.25em
\item 啟示錄
\begin{itemize}
\itemsep-.25em
\item 3:10 『你既遵守我忍耐的道，我必在普天下人受試煉的時候，保守你免去你的試煉。 11 我必快來！你要持守你所有的，免得人奪去你的\underline{冠冕}。
\item 6:1-2: 我看見羔羊揭開七印中第一印的時候，就聽見四活物中的一個活物聲音如雷說：「你來！」 2 我就觀看，見有一匹\underline{白馬}。騎在馬上的拿著弓，並有\underline{冠冕}賜給他，他便出來，勝了又要勝。
\item 19:11-21: 我觀看，見天開了。有一匹白馬，騎在馬上的稱為「誠信真實」，他審判、爭戰都按著公義。 12 他的眼睛如火焰，他頭上戴著許多冠冕，又有寫著的名字，除了他自己沒有人知道。 13 他穿著濺了血的衣服，他的名稱為「神之道」。 14 在天上的眾軍騎著白馬，穿著細麻衣，又白又潔，跟隨他。 15 有利劍從他口中出來，可以擊殺列國。他必用鐵杖轄管 b他們，並要踹全能神烈怒的酒榨。 16 在他衣服和大腿上有名寫著說：「萬王之王，萬主之主。」17 我又看見一位天使站在日頭中，向天空所飛的鳥大聲喊著說：「你們聚集來赴神的大筵席！ 18 可以吃君王與將軍的肉，壯士與馬和騎馬者的肉，並一切自主的、為奴的，以及大小人民的肉。」19 我看見那獸和地上的君王，並他們的眾軍都聚集，要與騎白馬的並他的軍兵爭戰。 20 那獸被擒拿，那在獸面前曾行奇事、迷惑受獸印記和拜獸像之人的假先知，也與獸同被擒拿。他們兩個就活活地被扔在燒著硫磺的火湖裡。 21 其餘的被騎白馬者口中出來的劍殺了。飛鳥都吃飽了他們的肉。
\end{itemize}
\item 撒迦利亞書
\begin{itemize}
\itemsep-.25em
\item 1:7-11 大流士第二年十一月，就是細罷特月二十四日，耶和華的話臨到易多的孫子、比利家的兒子先知撒迦利亞說： 8 我夜間觀看，見一人騎著紅馬，站在窪地番石榴樹中間，在他身後又有\underline{紅馬、黃馬和白馬}。 9 我對與我說話的天使說：「主啊，這是什麼意思？」他說：「我要指示你這是什麼意思。」 10 那站在番石榴樹中間的人說：\underline{「這是奉耶和華差遣在遍地走來走去的。」} 11 那些騎馬的對站在番石榴樹中間耶和華的使者說：「我們已在遍地走來走去，見全地都安息平靜。」
\item 6:1-8 我又舉目觀看，見有四輛車從兩山中間出來，那山是銅山。 2 第一輛車套著紅馬，第二輛車套著黑馬， 3 第三輛車套著白馬，第四輛車套著有斑點的壯馬。 4 我就問與我說話的天使說：「主啊，這是什麼意思？」 5 天使回答我說：「這是天的四風，是從普天下的主面前出來的。 6 套著黑馬的車往北方去，白馬跟隨在後，有斑點的馬往南方去。」 7 壯馬出來，要在遍地走來走去。天使說：「你們只管在遍地走來走去。」牠們就照樣行了。 8 他又呼叫我說：「看哪，\underline{往北方去的(黑馬白馬)已在北方安慰我的心。}」
\end{itemize}
\end{itemize}

\subsubsection{拿著弓(詩篇 78:9, 王下 13:17, 耶 50:14, 51:3,創49:22-24)}
\begin{itemize}
\itemsep-.25em
\item 詩篇 78:9: 以法蓮的子孫帶著兵器拿著弓，不要像他們的祖宗，是頑梗悖逆居心不正之輩，向著神心不誠實。以法蓮的子孫帶著兵器拿著弓，臨陣之日轉身退後。他們不遵守神的約，不肯
\item 列王紀下 13:17 說：「你開朝東的窗戶。」他就開了。以利沙說：「射箭吧。」他就射箭。以利沙說：「這是耶和華的得勝箭，就是戰勝亞蘭人的箭。因為你必在亞弗攻打亞蘭人，直到滅盡他們。」
\item 耶利米书
\begin{itemize}
\itemsep-.25em
\item 50:14 所有拉弓的，你們要在巴比倫的四圍擺陣，射箭攻擊她，不要愛惜箭枝，因她得罪了耶和華。
\item 51:3 我要打發外邦人來到巴比倫，簸揚她，使她的地空虛，在她遭禍的日子他們要周圍攻擊她。拉弓的要向拉弓的和貫甲挺身的射箭，不要憐惜她的少年人，要滅盡她的全軍。
\end{itemize}
\item 創 49:22-24: 22約瑟是多結果子的樹枝，是泉旁多結果的枝子，他的枝條探出牆外。弓箭手將他苦害，向他射箭，逼迫他。24 但他的弓仍舊堅硬，他的手健壯敏捷，這是因以色列的牧者，以色列的磐石，就是雅各的大能者。
\end{itemize}

\subsubsection{箭(賽 49:1-13, 撒迦利亞書 9:14 ,哈巴谷書 9:2-15)}
\begin{description}
\item 以賽亞書 49:1-13 \\
「眾海島啊，當聽我言！遠方的眾民哪，留心而聽！自我出胎，耶和華就選召我；自出母腹，他就提我的名。 2 他使我的口如快刀，將我藏在他手蔭之下；又使我成為磨亮的箭，將我藏在他箭袋之中。 3 對我說：『你是我的僕人以色列，我必因你得榮耀。』 4 我卻說：『我勞碌是徒然，我盡力是虛無虛空，然而我當得的理必在耶和華那裡，我的賞賜必在我神那裡。』5 「耶和華從我出胎造就我做他的僕人，要使雅各歸向他，使以色列到他那裡聚集。原來耶和華看我為尊貴，我的神也成為我的力量。 6 現在他說：『你做我的僕人，使雅各眾支派復興，使以色列中得保全的歸回，尚為小事，我還要使你做外邦人的光，叫你施行我的救恩，直到地極。』」 7 救贖主以色列的聖者耶和華，對那被人所藐視、本國所憎惡、官長所虐待的如此說：「君王要看見就站起，首領也要下拜，都因信實的耶和華，就是揀選你以色列的聖者。」8 耶和華如此說：「在悅納的時候，我應允了你；在拯救的日子，我濟助了你。我要保護你，使你做眾民的中保 a，復興遍地，使人承受荒涼之地為業。 9 對那被捆綁的人說：『出來吧！』對那在黑暗的人說：『顯露吧！』他們在路上必得飲食，在一切淨光的高處必有食物。 10 不飢不渴，炎熱和烈日必不傷害他們，因為憐恤他們的必引導他們，領他們到水泉旁邊。 11 我必使我的眾山成為大道，我的大路也被修高。 12 看哪，這些從遠方來，這些從北方，從西方來，這些從秦 b國來。」 13 諸天哪，應當歡呼！大地啊，應當快樂！眾山哪，應當發聲歌唱！因為耶和華已經安慰他的百姓，也要憐恤他困苦之民。
\item 撒迦利亞書 9:14 \\耶和華必顯現在他們以上，他的箭必射出像閃電。主耶和華必吹角，乘南方的旋風而行。
\item 哈巴谷書 9:2-15 \\耶和華啊，我聽見你的名聲就懼怕。耶和華啊，求你在這些年間復興你的作為，在這些年間顯明出來，在發怒的時候以憐憫為念。 3 神從提幔而來，聖者從巴蘭山臨到。（細拉）他的榮光遮蔽諸天，頌讚充滿大地。 4 他的輝煌如同日光，從他手裡射出光線，在其中藏著他的能力。 5 在他前面有瘟疫流行，在他腳下有熱症發出。 6 他站立，量了大地；觀看，趕散萬民。永久的山崩裂，長存的嶺塌陷，他的作為與古時一樣。 7 我見古珊的帳篷遭難，米甸的幔子戰兢。 8 耶和華啊，你乘在馬上，坐在得勝的車上，豈是不喜悅江河，向江河發怒氣，向洋海發憤恨嗎？ 9 你的弓全然顯露，向眾支派所起的誓都是可信的。（細拉）你以江河分開大地。 10 山嶺見你，無不戰懼；大水氾濫過去，深淵發聲，洶湧翻騰 c。 11 因你的箭射出發光，你的槍閃出光耀，日月都在本宮停住。 12 你發憤恨通行大地，發怒氣責打列國如同打糧。 13 你出來要拯救你的百姓，拯救你的受膏者，打破惡人家長的頭，露出他的腳 d，直到頸項。（細拉） 14 你用敵人的戈矛刺透他戰士的頭，他們來如旋風，要將我們分散，他們所喜愛的是暗中吞吃貧民。 15 你乘馬踐踏紅海，就是踐踏洶湧的大水。
\end{description}


\subsubsection{冠冕(林前9:25/啟 3:11/希 2:7/提後2:4/彼前5:4)}
\begin{itemize}
\itemsep-.25em
\item 哥林多前書 9:25 凡較力爭勝的，諸事都有節制。他們不過是要得能壞的冠冕，我們卻是要得不能壞的冠冕。
\item 提摩太後書 2:4 凡在軍中當兵的，不將世務纏身，好叫那招他當兵的人喜悅人若在場上比武，非按規矩，就不能得冠冕。
\item 啟示錄 3:11 我必快來！你要持守你所有的，免得人奪去你的冠冕。
\item 希伯來書 2:7 你叫他比天使微小一點，賜他榮耀、尊貴為冠冕，並將你手所造的都派他管理，
\item 彼得前書 5:4 到了牧長顯現的時候，你們必得那永不衰殘的榮耀冠冕。
\end{itemize}

\subsubsection{勝了又要勝 (啟13:1-6)}
我又看見一個獸從海中上來，有十角七頭，在十角上戴著十個冠冕，七頭上有褻瀆的名號。 2 我所看見的獸形狀像豹，腳像熊的腳，口像獅子的口。那龍將自己的能力、座位和大權柄都給了牠。 3 我看見獸的七頭中，有一個似乎受了死傷，那死傷卻醫好了。全地的人都稀奇，跟從那獸， 4 又拜那龍，因為牠將自己的權柄給了獸，也拜獸，說：「誰能比這獸，誰能與牠交戰呢？」 5 又賜給牠說誇大褻瀆話的口，又有權柄賜給牠，可以任意而行四十二個月。 6 獸就開口向神說褻瀆的話，褻瀆神的名並他的帳幕，以及那些住在天上的。又任憑牠與聖徒爭戰，並且得勝；也把權柄賜給牠，制伏各族、各民、各方、各國。


\subsection{錫安山}


\subsubsection{詩篇 2:6}
說：「我已經立我的君在錫安我的聖山上了。」
\subsubsection{希伯來書 12:18-29}
18 你們原不是來到那能摸的山，此山有火焰、密雲、黑暗、暴風、 19 角聲與說話的聲音。那些聽見這聲音的，都求不要再向他們說話， 20 因為他們當不起所命他們的話說：「靠近這山的，即便是走獸，也要用石頭打死。」 21 所見的極其可怕，甚至摩西說：「我甚是恐懼戰兢。」 22 你們乃是來到錫安山，永生神的城邑，就是天上的耶路撒冷；那裡有千萬的天使， 23 有名錄在天上諸長子之會所共聚的總會，有審判眾人的神和被成全之義人的靈魂， 24 並新約的中保耶穌以及所灑的血。這血所說的比亞伯的血所說的更美。 25 你們總要謹慎，不可棄絕那向你們說話的。因為那些棄絕在地上警戒他們的，尚且不能逃罪，何況我們違背那從天上警戒我們的呢？ 26 當時他的聲音震動了地，但如今他應許說：「再一次我不單要震動地，還要震動天。」 27 這再一次的話，是指明被震動的，就是受造之物都要挪去，使那不被震動的常存。 28 所以，我們既得了不能震動的國，就當感恩，照神所喜悅的，用虔誠、敬畏的心侍奉神， 29 因為我們的神乃是烈火。
\subsubsection{啟示錄 14:1-5 }
我又觀看，見羔羊站在錫安山，同他又有十四萬四千人，都有他的名和他父的名寫在額上。 2 我聽見從天上有聲音，像眾水的聲音和大雷的聲音，並且我所聽見的好像彈琴的所彈的琴聲。 3 他們在寶座前，並在四活物和眾長老前唱歌，彷彿是新歌，除了從地上買來的那十四萬四千人以外，沒有人能學這歌。 4 這些人未曾沾染婦女，他們原是童身。羔羊無論往哪裡去，他們都跟隨他。他們是從人間買來的，做初熟的果子歸於神和羔羊。 5 在他們口中察不出謊言來，他們是沒有瑕疵的。

\subsection{小书卷}

\subsubsection{以賽亞書 8:16，28:7，29:9-12，18}
\begin{itemize}
\itemsep-.25em
\item 8:16 你要捲起律法書，在我門徒中間封住訓誨。
\item 28:7 就是這地的人，也因酒搖搖晃晃，因濃酒東倒西歪。祭司和先知因濃酒搖搖晃晃，被酒所困，因濃酒東倒西歪。他們錯解默示，謬行審判。
\item 29:9 你們等候驚奇吧！你們宴樂昏迷吧！他們醉了，卻非因酒；他們東倒西歪，卻非因濃酒。 10 因為耶和華將沉睡的靈澆灌你們，封閉你們的眼，蒙蓋你們的頭。你們的眼就是先知，你們的頭就是先見。 11 所有的默示，你們看如封住的書卷。人將這書卷交給識字的，說：「請念吧。」他說：「我不能念，因為是封住了。」 12 又將這書卷交給不識字的人，說：「請念吧。」他說：「我不識字。」
\item 29:18 那時，聾子必聽見這書上的話，瞎子的眼必從迷蒙黑暗中得以看見。
\end{itemize}

\subsubsection{耶利米書 15:16/36:1-8}
\begin{itemize}
\itemsep-.25em
\item 耶和華萬軍之神啊，我得著你的言語就當食物吃了，你的言語是我心中的歡喜快樂，因我是稱為你名下的人。
\item 猶大王約西亞的兒子約雅敬第四年，耶和華的話臨到耶利米說： 2「你取一書卷，將我對你說攻擊以色列和猶大並各國的一切話，從我對你說話的那日，就是從約西亞的日子起，直到今日，都寫在其上。 3或者猶大家聽見我想要降於他們的一切災禍，各人就回頭離開惡道，我好赦免他們的罪孽和罪惡。所以耶利米召了尼利亞的兒子巴錄來，巴錄就從耶利米口中，將耶和華對耶利米所說的一切話寫在書卷上。 5 耶利米吩咐巴錄說：「我被拘管，不能進耶和華的殿， 6 所以你要去趁禁食的日子，在耶和華殿中將耶和華的話，就是你從我口中所寫在書卷上的話，念給百姓和一切從猶大城邑出來的人聽。 7 或者他們在耶和華面前懇求，各人回頭離開惡道，因為耶和華向這百姓所說要發的怒氣和憤怒是大的。」 8 尼利亞的兒子巴錄就照先知耶利米一切所吩咐的去行，在耶和華的殿中從書上念耶和華的話。
\end{itemize}

\subsubsection{以西結書 2:8-10，3:1-11}
\begin{itemize}
\itemsep-.25em
\item 「人子啊，要聽我對你所說的話，不要悖逆，像那悖逆之家。你要開口吃我所賜給你的。」 9 我觀看，見有一隻手向我伸出來，手中有一書卷。 10 他將書卷在我面前展開，內外都寫著字，其上所寫的有哀號、嘆息、悲痛的話。
\item 1 他對我說：「人子啊，要吃你所得的，要吃這書卷，好去對以色列家講說。」 2 於是我開口，他就使我吃這書卷。 3 又對我說：「人子啊，要吃我所賜給你的這書卷，充滿你的肚腹。」我就吃了，口中覺得其甜如蜜。
他對我說：「人子啊，要吃你所得的，要吃這書卷，好去對以色列家講說。」 2 於是我開口，他就使我吃這書卷。 3 又對我說：「人子啊，要吃我所賜給你的這書卷，充滿你的肚腹。」我就吃了，口中覺得其甜如蜜。4 他對我說：「人子啊，你往以色列家那裡去，將我的話對他們講說。 5 你奉差遣不是往那說話深奧、言語難懂的民那裡去，乃是往以色列家去。 6 不是往那說話深奧、言語難懂的多國去，他們的話語是你不懂得的，我若差你往他們那裡去，他們必聽從你。 7 以色列家卻不肯聽從你，因為他們不肯聽從我，原來以色列全家是額堅心硬的人。 8 看哪，我使你的臉硬過他們的臉，使你的額硬過他們的額。 9 我使你的額像金鋼鑽，比火石更硬。他們雖是悖逆之家，你不要怕他們，也不要因他們的臉色驚惶。」 10 他又對我說：「人子啊，我對你所說的一切話，要心裡領會，耳中聽聞。 11 你往你本國被擄的子民那裡去，他們或聽或不聽，你要對他們講說，告訴他們：『這是主耶和華說的。』」
\end{itemize}

\subsubsection{但以理書 12:4-13}
但以理啊，你要隱藏這話，封閉這書，直到末時。必有多人來往奔跑，知識就必增長。」5 我但以理觀看，見另有兩個人站立，一個在河這邊，一個在河那邊。 6 有一個問那站在河水以上穿細麻衣的說：「這奇異的事到幾時才應驗呢？」 7 我聽見那站在河水以上穿細麻衣的，向天舉起左右手，指著活到永遠的主起誓說：「要到一載，二載，半載，打破聖民權力的時候，這一切事就都應驗了。」 8 我聽見這話，卻不明白，就說：「我主啊，這些事的結局是怎樣呢？」 9 他說：「但以理啊，你只管去，因為這話已經隱藏封閉，直到末時。 10 必有許多人使自己清淨潔白，且被熬煉，但惡人仍必行惡。一切惡人都不明白，唯獨智慧人能明白。 11 從除掉常獻的燔祭，並設立那行毀壞可憎之物的時候，必有一千二百九十日。 12 等到一千三百三十五日的，那人便為有福！ 13 你且去等候結局，因為你必安歇。到了末期，你必起來，享受你的福分。」

\subsubsection{撒迦利亞書 5:1-4}
我又舉目觀看，見有一飛行的書卷。 2 他問我說：「你看見什麼？」我回答說：「我看見一飛行的書卷，長二十肘，寬十肘。」 3 他對我說：「這是發出行在遍地上的咒詛，凡偷竊的必按卷上這面的話除滅，凡起假誓的必按卷上那面的話除滅。 4 萬軍之耶和華說：我必使這書卷出去，進入偷竊人的家和指我名起假誓人的家，必常在他家裡，連房屋帶木石都毀滅了。」



\subsubsection{啟 5:1-10，11:1-11}
\begin{itemize} 
\itemsep-.25em
\item 5:1 我看見坐寶座的右手中有書卷，裡外都寫著字，用七印封嚴了。 2 我又看見一位大力的天使大聲宣傳說：「有誰配展開那書卷，揭開那七印呢？」 3 在天上、地上、地底下，沒有能展開、能觀看那書卷的。 4 因為沒有配展開、配觀看那書卷的，我就大哭。 5 長老中有一位對我說：「不要哭。看哪，猶大支派中的獅子，大衛的根，他已得勝，能以展開那書卷、揭開那七印！」6 我又看見寶座與四活物並長老之中有羔羊站立，像是被殺過的，有七角七眼，就是神的七靈，奉差遣往普天下去的。 7 這羔羊前來，從坐寶座的右手裡拿了書卷。 8 他既拿了書卷，四活物和二十四位長老就俯伏在羔羊面前，各拿著琴和盛滿了香的金爐。這香就是眾聖徒的祈禱。 9 他們唱新歌，說：「你配拿書卷，配揭開七印，因為你曾被殺，用自己的血從各族、各方、各民、各國中買了人來，叫他們歸於神， 10 又叫他們成為國民做祭司，歸於神，在地上執掌王權。」
\item 我又看見另有一位大力的天使從天降下，披著雲彩，頭上有虹，臉面像日頭，兩腳像火柱。 2 他手裡拿著小書卷，是展開的。他右腳踏海，左腳踏地， 3 大聲呼喊，好像獅子吼叫。呼喊完了，就有七雷發聲。 4 七雷發聲之後，我正要寫出來，就聽見從天上有聲音說：「七雷所說的你要封上，不可寫出來。」 5 我所看見的那踏海踏地的天使向天舉起右手來， 6 指著那創造天和天上之物、地和地上之物、海和海中之物，直活到永永遠遠的，起誓說：「不再有時日了！ 7 但在第七位天使吹號發聲的時候，神的奧祕就成全了，正如神所傳給他僕人眾先知的佳音。」8 我先前從天上所聽見的那聲音又吩咐我說：「你去，把那踏海踏地之天使手中展開的小書卷取過來。」 9 我就去到天使那裡，對他說：「請你把小書卷給我。」他對我說：「你拿著吃盡了，便叫你肚子發苦，然而在你口中要甜如蜜。」 10 我從天使手中把小書卷接過來，吃盡了，在我口中果然甜如蜜；吃了以後，肚子覺得發苦了。 11 天使對我說：「你必指著多民、多國、多方、多王再說預言。」
\end{itemize} 


\subsection{兩刃利劍——別迦摩}

\subsubsection{以賽亞書 11:1-5, 49:1-13}
\begin{itemize} 
\itemsep-.25em
\item 從耶西的本必發一條，從他根生的枝子必結果實。 2 耶和華的靈必住在他身上，就是使他有智慧和聰明的靈，謀略和能力的靈，知識和敬畏耶和華的靈。 3 他必以敬畏耶和華為樂，行審判不憑眼見，斷是非也不憑耳聞； 4 卻要以公義審判貧窮人，以正直判斷世上的謙卑人，以口中的杖擊打世界，以嘴裡的氣殺戮惡人。 5 公義必當他的腰帶，信實必當他脅下的帶子。
\item 1 「眾海島啊，當聽我言！遠方的眾民哪，留心而聽！自我出胎，耶和華就選召我；自出母腹，他就提我的名。 2 他使我的口如快刀，將我藏在他手蔭之下；又使我成為磨亮的箭，將我藏在他箭袋之中。 3 對我說：『你是我的僕人以色列，我必因你得榮耀。』 4 我卻說：『我勞碌是徒然，我盡力是虛無虛空，然而我當得的理必在耶和華那裡，我的賞賜必在我神那裡。』5 「耶和華從我出胎造就我做他的僕人，要使雅各歸向他，使以色列到他那裡聚集。原來耶和華看我為尊貴，我的神也成為我的力量。 6 現在他說：『你做我的僕人，使雅各眾支派復興，使以色列中得保全的歸回，尚為小事，我還要使你做外邦人的光，叫你施行我的救恩，直到地極。』」 7 救贖主以色列的聖者耶和華，對那被人所藐視、本國所憎惡、官長所虐待的如此說：「君王要看見就站起，首領也要下拜，都因信實的耶和華，就是揀選你以色列的聖者。」8 耶和華如此說：「在悅納的時候，我應允了你；在拯救的日子，我濟助了你。我要保護你，使你做眾民的中保，復興遍地，使人承受荒涼之地為業。 9 對那被捆綁的人說：『出來吧！』對那在黑暗的人說：『顯露吧！』他們在路上必得飲食，在一切淨光的高處必有食物。 10 不飢不渴，炎熱和烈日必不傷害他們，因為憐恤他們的必引導他們，領他們到水泉旁邊。 11 我必使我的眾山成為大道，我的大路也被修高。 12 看哪，這些從遠方來，這些從北方，從西方來，這些從秦國來。」 13 諸天哪，應當歡呼！大地啊，應當快樂！眾山哪，應當發聲歌唱！因為耶和華已經安慰他的百姓，也要憐恤他困苦之民。
\end{itemize}

\subsubsection{以弗所書 6:17}
並戴上救恩的頭盔，拿著聖靈的寶劍，就是神的道。

\subsubsection{希伯來書 4:12-13}
12 神的道是活潑的，是有功效的，比一切兩刃的劍更快，甚至魂與靈、骨節與骨髓，都能刺入、剖開，連心中的思念和主意都能辨明； 13 並且被造的沒有一樣在他面前不顯然的，原來萬物在那與我們有關係的主眼前，都是赤露敞開的


\subsubsection{啟示錄 1:12-16, 2:12-17, 19:11-16}
\begin{itemize} 
\itemsep-.25em
\item 12 我轉過身來，要看是誰發聲與我說話；既轉過來，就看見七個金燈臺， 13 燈臺中間有一位好像人子，身穿長衣，直垂到腳，胸間束著金帶。 14 他的頭與髮皆白，如白羊毛，如雪，眼目如同火焰， 15 腳好像在爐中鍛煉光明的銅，聲音如同眾水的聲音。 16 他右手拿著七星，從他口中出來一把兩刃的利劍，面貌如同烈日放光。
\item 「你要寫信給別迦摩教會的使者說：『那有兩刃利劍的說： 13 我知道你的居所，就是有撒旦座位之處。當我忠心的見證人安提帕在你們中間、撒旦所住的地方被殺之時，你還堅守我的名，沒有棄絕我的道。 14 然而，有幾件事我要責備你，因為在你那裡，有人服從了巴蘭的教訓。這巴蘭曾教導巴勒將絆腳石放在以色列人面前，叫他們吃祭偶像之物，行姦淫的事。 15 你那裡也有人照樣服從了尼哥拉一黨人的教訓。『所以，你當悔改！若不悔改，我就快臨到你那裡，用我口中的劍攻擊他們。 17 聖靈向眾教會所說的話，凡有耳的，就應當聽！得勝的，我必將那隱藏的嗎哪賜給他，並賜他一塊白石，石上寫著新名，除了那領受的以外，沒有人能認識。』
\item 我觀看，見天開了。有一匹白馬，騎在馬上的稱為「誠信真實」，他審判、爭戰都按著公義。 12 他的眼睛如火焰，他頭上戴著許多冠冕，又有寫著的名字，除了他自己沒有人知道。 13 他穿著濺了血的衣服，他的名稱為「神之道」。 14 在天上的眾軍騎著白馬，穿著細麻衣，又白又潔，跟隨他。 15 有利劍從他口中出來，可以擊殺列國。他必用鐵杖轄管 b他們，並要踹全能神烈怒的酒榨。 16 在他衣服和大腿上有名寫著說：「萬王之王，萬主之主。」
\end{itemize} 



\subsection{隱藏的嗎哪——別迦摩}

\subsubsection{出埃及記 16:4-36}
耶和華對摩西說：「我要將糧食從天降給你們。百姓可以出去，每天收每天的份，我好試驗他們遵不遵我的法度。 5 到第六天，他們要把所收進來的預備好了，比每天所收的多一倍。」 6 摩西、亞倫對以色列眾人說：「到了晚上，你們要知道是耶和華將你們從埃及地領出來的。 7 早晨，你們要看見耶和華的榮耀，因為耶和華聽見你們向他所發的怨言了。我們算什麼，你們竟向我們發怨言呢？」 8 摩西又說：「耶和華晚上必給你們肉吃，早晨必給你們食物得飽，因為你們向耶和華發的怨言，他都聽見了。我們算什麼？你們的怨言不是向我們發的，乃是向耶和華發的。」 9 摩西對亞倫說：「你告訴以色列全會眾說：『你們就近耶和華面前，因為他已經聽見你們的怨言了。』」 10 亞倫正對以色列全會眾說話的時候，他們向曠野觀看，不料，耶和華的榮光在雲中顯現。 11 耶和華曉諭摩西說： 12 「我已經聽見以色列人的怨言。你告訴他們說：『到黃昏的時候，你們要吃肉，早晨必有食物得飽，你們就知道我是耶和華你們的神。』」13 到了晚上，有鵪鶉飛來，遮滿了營。早晨，在營四圍的地上有露水。 14 露水上升之後，不料，野地面上有如白霜的小圓物。 15 以色列人看見，不知道是什麼，就彼此對問說：「這是什麼呢？」摩西對他們說：「這就是耶和華給你們吃的食物。 16 耶和華所吩咐的是這樣：你們要按著各人的飯量，為帳篷裡的人按著人數收起來，各拿一俄梅珥。」 17 以色列人就這樣行，有多收的，有少收的。 18 及至用俄梅珥量一量，多收的也沒有餘，少收的也沒有缺，各人按著自己的飯量收取。 19 摩西對他們說：「所收的，不許什麼人留到早晨。」 20 然而他們不聽摩西的話，內中有留到早晨的，就生蟲變臭了。摩西便向他們發怒。21 他們每日早晨，按著各人的飯量收取，日頭一發熱，就消化了。 22 到第六天，他們收了雙倍的食物，每人兩俄梅珥。會眾的官長來告訴摩西， 23 摩西對他們說：「耶和華這樣說：『明天是聖安息日，是向耶和華守的聖安息日。你們要烤的就烤了，要煮的就煮了，所剩下的都留到早晨。』」 24 他們就照摩西的吩咐留到早晨，也不臭，裡頭也沒有蟲子。 25 摩西說：「你們今天吃這個吧，因為今天是向耶和華守的安息日，你們在田野必找不著了。 26 六天可以收取，第七天乃是安息日，那一天必沒有了。」 27 第七天，百姓中有人出去收，什麼也找不著。 28 耶和華對摩西說：「你們不肯守我的誡命和律法，要到幾時呢？ 29 你們看，耶和華既將安息日賜給你們，所以第六天他賜給你們兩天的食物，第七天各人要住在自己的地方，不許什麼人出去。」 30 於是百姓第七天安息了。31 這食物，以色列家叫嗎哪，樣子像芫荽子，顏色是白的，滋味如同攙蜜的薄餅。 32 摩西說：「耶和華所吩咐的是這樣：要將一滿俄梅珥嗎哪留到世世代代，使後人可以看見我當日將你們領出埃及地，在曠野所給你們吃的食物。」 33 摩西對亞倫說：「你拿一個罐子，盛一滿俄梅珥嗎哪，存在耶和華面前，要留到世世代代。」 34 耶和華怎麼吩咐摩西，亞倫就怎麼行，把嗎哪放在法櫃前存留。 35 以色列人吃嗎哪共四十年，直到進了有人居住之地，就是迦南的境界。 36 （俄梅珥乃伊法十分之一。）

\subsubsection{民數記 11:4-9}
4 他們中間的閒雜人大起貪慾的心。以色列人又哭號說：「誰給我們肉吃呢？ 5 我們記得在埃及的時候不花錢就吃魚，也記得有黃瓜、西瓜、韭菜、蔥、蒜。 6 現在我們的心血枯竭了，除這嗎哪以外，在我們眼前並沒有別的東西！」 7 這嗎哪彷彿芫荽子，又好像珍珠。 8 百姓周圍行走，把嗎哪收起來，或用磨推，或用臼搗，煮在鍋中，又做成餅，滋味好像新油。 9 夜間露水降在營中，嗎哪也隨著降下。


\subsubsection{申命記 8:3}
他苦煉你，任你飢餓，將你和你列祖所不認識的嗎哪賜給你吃，使你知道人活著不是單靠食物，乃是靠耶和華口裡所出的一切話。

\subsubsection{約書亞記 5:12}
他們吃了那地的出產，第二日嗎哪就止住了，以色列人也不再有嗎哪了。那一年，他們卻吃迦南地的出產。

\subsubsection{尼希米記 9:15,20}
15從天上賜下糧食充他們的飢，從磐石使水流出解他們的渴，又吩咐他們進去得你起誓應許賜給他們的地。20 你也賜下你良善的靈教訓他們，未嘗不賜嗎哪使他們糊口，並賜水解他們的渴。

\subsubsection{詩篇 78:18-25}
他們心中試探神，隨自己所欲的求食物，
19 並且妄論神說：「神在曠野豈能擺設筵席嗎？
20 他曾擊打磐石，使水湧出成了江河，他還能賜糧食嗎？還能為他的百姓預備肉嗎？」
21 所以耶和華聽見就發怒，有烈火向雅各燒起，有怒氣向以色列上騰。
22 因為他們不信服神，不倚賴他的救恩。
23 他卻吩咐天空，又敞開天上的門，
24 降嗎哪像雨給他們吃，將天上的糧食賜給他們。
25 各人吃大能者的食物，他賜下糧食，使他們飽足。


\subsubsection{約翰福音 6:30-33，47-58}
30 他們又說：「你行什麼神蹟，叫我們看見就信你？你到底做什麼事呢？ 31 我們的祖宗在曠野吃過嗎哪，如經上寫著說：『他從天上賜下糧來給他們吃。』」 32 耶穌說：「我實實在在地告訴你們：那從天上來的糧不是摩西賜給你們的，乃是我父將天上來的真糧賜給你們。 33 因為神的糧就是那從天上降下來賜生命給世界的。」 47 我實實在在地告訴你們：信的人有永生。 48 我就是生命的糧。 49 你們的祖宗在曠野吃過嗎哪，還是死了； 50 這是從天上降下來的糧，叫人吃了就不死。 51 我是從天上降下來生命的糧，人若吃這糧，就必永遠活著。我所要賜的糧就是我的肉，為世人之生命所賜的。」52 因此，猶太人彼此爭論說：「這個人怎能把他的肉給我們吃呢？」 53 耶穌說：「我實實在在地告訴你們：你們若不吃人子的肉，不喝人子的血，就沒有生命在你們裡面。 54 吃我肉、喝我血的人就有永生，在末日我要叫他復活。 55 我的肉真是可吃的，我的血真是可喝的。 56 吃我肉、喝我血的人常在我裡面，我也常在他裡面。 57 永活的父怎樣差我來，我又因父活著，照樣，吃我肉的人也要因我活著。 58 這就是從天上降下來的糧。吃這糧的人就永遠活著，不像你們的祖宗吃過嗎哪還是死了。」

\subsubsection{哥林多前書 10:3，4}
並且都吃了一樣的靈食， 4也都喝了一樣的靈水；所喝的是出於隨著他們的靈磐石，那磐石就是基督。

\subsubsection{希伯來書 9:4}
有金香爐，有包金的約櫃，櫃裡有盛嗎哪的金罐和亞倫發過芽的杖並兩塊約版，

\subsubsection{啟示錄 2:17}
聖靈向眾教會所說的話，凡有耳的，就應當聽！得勝的，我必將那隱藏的嗎哪賜給他，並賜他一塊白石，石上寫著新名，除了那領受的以外，沒有人能認識。

\subsection{一塊白石——別迦摩}

\subsubsection{神改的名字-亞伯拉罕,撒拉,雅各}
\begin{itemize} 
\itemsep-.25em
\item 創世記 17：1-8
亞伯蘭年九十九歲的時候，耶和華向他顯現，對他說：「我是全能的神，你當在我面前做完全人。 2 我就與你立約，使你的後裔極其繁多。」 3 亞伯蘭俯伏在地，神又對他說： 4 「我與你立約：你要做多國的父。 5 從此以後，你的名不再叫亞伯蘭，要叫亞伯拉罕，因為我已立你做多國的父。 6 我必使你的後裔極其繁多，國度從你而立，君王從你而出。 7 我要與你並你世世代代的後裔堅立我的約，做永遠的約，是要做你和你後裔的神。 8 我要將你現在寄居的地，就是迦南全地，賜給你和你的後裔永遠為業，我也必做他們的神。」
\item 創世記 17：15-16，19
神又對亞伯拉罕說：「你的妻子撒萊，不可再叫撒萊，她的名要叫撒拉。 16 我必賜福給她，也要使你從她得一個兒子。我要賜福給她，她也要做多國之母，必有百姓的君王從她而出。」 19神說：「不然，你妻子撒拉要給你生一個兒子，你要給他起名叫以撒。我要與他堅定所立的約，做他後裔永遠的約。
\item 創世記 32:27-29
…27那人說：「你名叫什麼？」他說：「我名叫雅各。」 28那人說：「你的名不要再叫雅各，要叫以色列，因為你與神與人較力，都得了勝。」 29雅各問他說：「請將你的名告訴我。」那人說：「何必問我的名？」於是在那裡給雅各祝福。
\item 創世記 35:9-11
雅各從巴旦亞蘭回來，神又向他顯現，賜福於他， 10且對他說：「你的名原是雅各，從今以後不要再叫雅各，要叫以色列。」這樣，他就改名叫以色列。 11神又對他說：「我是全能的神。你要生養眾多，將來有一族和多國的民從你而生，又有君王從你而出。
\item 列王紀上 18:31
以利亞照雅各子孫支派的數目，取了十二塊石頭（耶和華的話曾臨到雅各說「你的名要叫以色列」），
\end{itemize} 

\subsubsection{神赐的名字-亞當,何西阿的三个孩子，耶底底亞，以賽亞儿子}
\begin{itemize} 
\itemsep-.25em
\item 亞當的名字是神给取的。
\item 撒母耳記下 12：24-25
大衛安慰他的妻拔示巴，與她同寢，她就生了兒子，給他起名叫所羅門。耶和華也喜愛他，25就藉先知拿單賜他一個名字，叫耶底底亞，因為耶和華愛他
\item 以賽亞書 8:1-4
耶和華對我說：「你取一塊大板子，拿人的筆，寫上『瑪黑珥‧沙拉勒‧哈施‧罷斯』。 2 我要用可靠的證人，烏利亞祭司和耶比利家的兒子撒迦利亞為我作證。」「瑪黑珥‧沙拉勒‧哈施‧罷斯」意思是「擄掠速臨、搶奪即將發生」。3 我親近女先知；她就懷孕生子，耶和華對我說：「給他起名叫瑪黑珥‧沙拉勒‧哈施‧罷斯； 4 因為在這孩子還不曉得叫爸爸媽媽以前，大馬士革的財寶和撒瑪利亞的擄物必被亞述王掠奪一空。」
\item 何西阿书 1：4-9
耶和华就对何西阿说：“给他起名叫耶斯列；因为不久我就要惩罚耶户家在耶斯列流人血的罪，我要灭绝以色列家的国。5 “到那日，我要在耶斯列平原折断以色列的弓。”6 后来，歌玛又怀孕，生了一个女儿。耶和华对何西阿说：“给她起名叫罗．路哈玛（“罗．路哈玛”意即“不蒙怜恤”）；因为我必不再怜恤以色列家，我决不赦免他们。7 “但我要怜恤犹大家，使他们靠耶和华他们的　神得救。我要拯救他们，并不是靠弓、靠刀、靠战争，也不是靠战马、靠骑兵。”8 罗．路哈玛断奶之后，歌玛又怀孕生了一个儿子。9 耶和华说：“给他起名叫罗．阿米（“罗．阿米”意即“非我民”）；因为你们不是我的子民，我也不是你们的神。
\end{itemize}

\subsubsection{新约中神赐的名字人-耶穌,磯法,约翰}
\begin{itemize} 
\itemsep-.25em
\item 馬太福音 1：18-25
18 耶穌基督降生的事記在下面：他母親馬利亞已經許配了約瑟，還沒有迎娶，馬利亞就從聖靈懷了孕。 19 她丈夫約瑟是個義人，不願意明明地羞辱她，想要暗暗地把她休了。 20 正思念這事的時候，有主的使者向他夢中顯現，說：「大衛的子孫約瑟，不要怕，只管娶過你的妻子馬利亞來，因她所懷的孕是從聖靈來的。 21 她將要生一個兒子，你要給他起名叫耶穌，因他要將自己的百姓從罪惡裡救出來。」 22 這一切的事成就，是要應驗主藉先知所說的話說： 23 「必有童女懷孕生子，人要稱他的名為以馬內利。」（「以馬內利」翻出來就是「神與我們同在」。） 24 約瑟醒了，起來，就遵著主使者的吩咐把妻子娶過來， 25 只是沒有和她同房，等她生了兒子，就給他起名叫耶穌。
\item 約翰福音 1： 42 於是領他去見耶穌。耶穌看著他說：「你是約翰的兒子西門，你要稱為磯法。」（「磯法」翻出來就是「彼得」。）
\item 路加福音 1:8-23 8 撒迦利亞按班次在神面前供祭司的職分， 9 照祭司的規矩掣籤，得進主殿燒香。 10 燒香的時候，眾百姓在外面禱告。 11 有主的使者站在香壇的右邊向他顯現。 12 撒迦利亞看見，就驚慌害怕。 13 天使對他說：「撒迦利亞，不要害怕，因為你的祈禱已經被聽見了。你的妻子伊利莎白要給你生一個兒子，你要給他起名叫約翰。 14 你必歡喜快樂，有許多人因他出世也必喜樂。 15 他在主面前將要為大，淡酒濃酒都不喝，從母腹裡就被聖靈充滿了。 16 他要使許多以色列人回轉，歸於主他們的神。 17 他必有以利亞的心志能力，行在主的前面，叫為父的心轉向兒女，叫悖逆的人轉從義人的智慧，又為主預備合用的百姓。」 18 撒迦利亞對天使說：「我憑著什麼可知道這事呢？我已經老了，我的妻子也年紀老邁了。」 19 天使回答說：「我是站在神面前的加百列，奉差而來對你說話，將這好信息報給你。 20 到了時候，這話必然應驗。只因你不信，你必啞巴，不能說話，直到這事成就的日子。」 21 百姓等候撒迦利亞，詫異他許久在殿裡。 22 及至他出來，不能和他們說話，他們就知道他在殿裡見了異象，因為他直向他們打手勢，竟成了啞巴。 23 他供職的日子已滿，就回家去了。
\end{itemize} 

\subsubsection{名字}
\begin{itemize} 
\itemsep-.25em
\item 以賽亞書 1:26
我也必復還你的審判官，像起初一樣；復還你的謀士，像起先一般。然後你必稱為公義之城，忠信之邑。」
\item 以賽亞書 43:1-7
1 雅各啊，創造你的耶和華，以色列啊，造成你的那位，現在如此說：「你不要害怕，因為我救贖了你；我曾提你的名召你，你是屬我的。 2 你從水中經過，我必與你同在；你趟過江河，水必不漫過你；你從火中行過，必不被燒，火焰也不著在你身上。 3 因為我是耶和華你的神，是以色列的聖者，你的救主。我已經使埃及做你的贖價，使古實和西巴代替你。 4 因我看你為寶為尊，又因我愛你，所以我使人代替你，使列邦人替換你的生命。 5 不要害怕，因我與你同在。我必領你的後裔從東方來，又從西方招聚你。 6 我要對北方說『交出來！』，對南方說『不要拘留！』，將我的眾子從遠方帶來，將我的眾女從地極領回， 7 就是凡稱為我名下的人，是我為自己的榮耀創造的，是我所做成、所造做的。」
\item 以賽亞書 48:19
你的後裔也必多如海沙，你腹中所生的也必多如沙粒，他的名在我面前必不剪除，也不滅絕。」
\item 以賽亞書 55:13
松樹長出代替荊棘，番石榴長出代替蒺藜，這要為耶和華留名，作為永遠的證據，不能剪除。」
\item 以賽亞書 56：1-5
1 耶和華如此說：「你們當守公平，行公義，因我的救恩臨近，我的公義將要顯現。 2 謹守安息日而不干犯，禁止己手而不作惡，如此行、如此持守的人，便為有福！」 3 與耶和華聯合的外邦人不要說：「耶和華必定將我從他民中分別出來。」太監也不要說：「我是枯樹。」 4 因為耶和華如此說：「那些謹守我的安息日，揀選我所喜悅的事，持守我約的太監， 5 我必使他們在我殿中、在我牆內有紀念，有名號，比有兒女的更美。我必賜他們永遠的名，不能剪除。
\item 以賽亞書 60:18
你地上不再聽見強暴的事，境內不再聽見荒涼毀滅的事。你必稱你的牆為『拯救』，稱你的門為『讚美』。
\item 以賽亞書 62：1-5
1我因錫安必不靜默，為耶路撒冷必不息聲，直到他的公義如光輝發出，他的救恩如明燈發亮。 2列國必見你的公義，列王必見你的榮耀。你必得新名的稱呼，是耶和華親口所起的。 你在耶和華的手中要作為華冠，在你神的掌上必作為冕旒。 4 你必不再稱為「撇棄的」，你的地也不再稱為「荒涼的」，你卻要稱為「我所喜悅的」，你的地也必稱為「有夫之婦」，因為耶和華喜悅你，你的地也必歸他。 5 少年人怎樣娶處女，你的眾民 a也要照樣娶你；新郎怎樣喜悅新婦，你的神也要照樣喜悅你。
\item 以賽亞書 66:22
耶和華說：「我所要造的新天新地怎樣在我面前長存，你們的後裔和你們的名字也必照樣長存。
\item 西番雅書 3:20
那時，我必領你們進來，聚集你們。我使你們被擄之人歸回的時候，就必使你們在地上的萬民中有名聲，得稱讚。」這是耶和華說的。
\item 以弗所書 3：14-19
因此，我在父面前屈膝—— 15 天上地上的各家，都是從他得名—— 16 求他按著他豐盛的榮耀，藉著他的靈，叫你們心裡的力量剛強起來， 17 使基督因你們的信住在你們心裡，叫你們的愛心有根有基， 18 能以和眾聖徒一同明白基督的愛是何等長闊高深， 19 並知道這愛是過於人所能測度的，便叫神一切所充滿的充滿了你們。
\item 啟示錄 2:17
聖靈向眾教會所說的話，凡有耳的，就應當聽！得勝的，我必將那隱藏的嗎哪賜給他，並賜他一塊白石，石上寫著新名，除了那領受的以外，沒有人能認識。』
\item 啟示錄 3:12
得勝的，我要叫他在我神殿中做柱子，他也必不再從那裡出去。我又要將我神的名和我神城的名——這城就是從天上、從我神那裡降下來的新耶路撒冷——並我的新名，都寫在他上面。
\item 啟示錄 19:12
他的眼睛如火焰，他頭上戴著許多冠冕，又有寫著的名字，除了他自己沒有人知道。
\item 啟示錄 19:16
有利劍從他口中出來，可以擊殺列國。他必用鐵杖轄管他們，並要踹全能神烈怒的酒榨。 16在他衣服和大腿上有名寫著說：「萬王之王，萬主之主。」
\end{itemize} 

\subsection{鐵杖權柄制伏列國——推雅推喇}
\subsubsection{詩篇 2:6-12，110:2}
\begin{itemize} 
\itemsep-.25em
\item 你求我，我就將列國賜你為基業，將地極賜你為田產。
7 受膏者說：「我要傳聖旨。耶和華曾對我說：『你是我的兒子，我今日生你。
8 你求我，我就將列國賜你為基業，將地極賜你為田產。
9 你必用鐵杖打破他們，你必將他們如同窯匠的瓦器摔碎。』」
10 現在你們君王應當醒悟，你們世上的審判官該受管教！
11 當存畏懼侍奉耶和華，又當存戰兢而快樂。
12 當以嘴親子，恐怕他發怒，你們便在道中滅亡，因為他的怒氣快要發作。凡投靠他的，都是有福的！
\item 耶和華必使你從錫安伸出能力的杖來，你要在你仇敵中掌權。
\end{itemize} 

\subsubsection{以賽亞書 11:1-6}
1 從耶西的本必發一條，從他根生的枝子必結果實。 2 耶和華的靈必住在他身上，就是使他有智慧和聰明的靈，謀略和能力的靈，知識和敬畏耶和華的靈。 3 他必以敬畏耶和華為樂，行審判不憑眼見，斷是非也不憑耳聞； 4 卻要以公義審判貧窮人，以正直判斷世上的謙卑人，以口中的杖擊打世界，以嘴裡的氣殺戮惡人。 5 公義必當他的腰帶，信實必當他脅下的帶子。

\subsubsection{耶利米書 19:11}
對他們說：『萬軍之耶和華如此說：我要照樣打碎這民和這城，正如人打碎窯匠的瓦器，以致不能再囫圇。並且人要在陀斐特葬埋屍首，甚至無處可葬。

\subsubsection{但以理書 2:44}
當那列王在位的時候，天上的神必另立一國，永不敗壞，也不歸別國的人，卻要打碎滅絕那一切國，這國必存到永遠。

\subsubsection{使徒行傳 15:28}
因為聖靈和我們定意不將別的重擔放在你們身上，唯有幾件事是不可少的，

\subsubsection{啟示錄 2:24-29，3:21，12:5，19:15，20:4 }
\begin{itemize} 
\itemsep-.25em
\item 2:24 『至於你們推雅推喇其餘的人，就是一切不從那教訓、不曉得他們素常所說撒旦深奧之理的人，我告訴你們：我不將別的擔子放在你們身上； 25 但你們已經有的，總要持守，直等到我來。 26 那得勝又遵守我命令到底的，我要賜給他權柄制伏列國， 27 他必用鐵杖轄管 b他們，將他們如同窯戶的瓦器打得粉碎，像我從我父領受的權柄一樣。 28 我又要把晨星賜給他。 29 聖靈向眾教會所說的話，凡有耳的，就應當聽！』
\item 得勝的，我要賜他在我寶座上與我同坐，就如我得了勝，在我父的寶座上與他同坐一般。
\item 3:21 婦人生了一個男孩子，是將來要用鐵杖轄管萬國的；她的孩子被提到神寶座那裡去了。
\item 19:15 有利劍從他口中出來，可以擊殺列國。他必用鐵杖轄管他們，並要踹全能神烈怒的酒榨。
\item 20:4 我又看見幾個寶座，也有坐在上面的，並有審判的權柄賜給他們。我又看見那些因為給耶穌作見證並為神之道被斬者的靈魂，和那沒有拜過獸與獸像，也沒有在額上和手上受過牠印記之人的靈魂；他們都復活了，與基督一同做王一千年。
\end{itemize} 

\subsection{晨星——推雅推喇}

\subsubsection{民數記 24:17}
我看他卻不在現時，我望他卻不在近日。有星要出於雅各，有杖要興於以色列，必打破摩押的四角，毀壞擾亂之子。

\subsubsection{約伯記 38：1-15}
1 那時，耶和華從旋風中回答約伯說： 2 「誰用無知的言語，使我的旨意暗昧不明？ 3 你要如勇士束腰，我問你，你可以指示我。 4 我立大地根基的時候，你在哪裡呢？你若有聰明，只管說吧。 5 你若曉得就說，是誰定地的尺度？是誰把準繩拉在其上？ 6 地的根基安置在何處？地的角石是誰安放的？ 7 那時晨星一同歌唱，神的眾子也都歡呼。8 「海水衝出，如出胎胞，那時誰將它關閉呢？ 9 是我用雲彩當海的衣服，用幽暗當包裹它的布， 10 為它定界限，又安門和閂， 11 說：『你只可到這裡，不可越過，你狂傲的浪要到此止住。』12 「你自生以來，曾命定晨光，使清晨的日光知道本位， 13 叫這光普照地的四極，將惡人從其中驅逐出來嗎？ 14 因這光，地面改變如泥上印印，萬物出現如衣服一樣。 15 亮光不照惡人，強橫的膀臂也必折斷。

\subsubsection{以賽亞書 8:19-21，14:12-20}
\begin{itemize} 
\itemsep-.25em
\item 19有人對你們說：「當求問那些交鬼的和行巫術的，就是聲音綿蠻、言語微細的。」你們便回答說：「百姓不當求問自己的神嗎？豈可為活人求問死人呢？ 20人當以訓誨和法度為標準。」他們所說的若不與此相符，必不得見晨光。 21他們必經過這地，受艱難，受飢餓，飢餓的時候心中焦躁，咒罵自己的君王和自己的神。…
\item 14:12 明亮之星，早晨之子啊，你何竟從天墜落？你這攻敗列國的，何竟被砍倒在地上？ 13 你心裡曾說：「我要升到天上，我要高舉我的寶座在神眾星以上，我要坐在聚會的山上，在北方的極處， 14 我要升到高雲之上，我要與至上者同等。」 15 然而你必墜落陰間，到坑中極深之處。 16 凡看見你的都要定睛看你，留意看你，說：「使大地戰抖，使列國震動， 17 使世界如同荒野，使城邑傾覆，不釋放被擄的人歸家，是這個人嗎？」 18 列國的君王俱各在自己陰宅的榮耀中安睡， 19 唯獨你被拋棄，不得入你的墳墓，好像可憎的枝子，以被殺的人為衣，就是被刀刺透，墜落坑中石頭那裡的。你又像被踐踏的屍首一樣。 20 你不得與君王同葬，因為你敗壞你的國，殺戮你的民。惡人後裔的名必永不提說！
\end{itemize} 

\subsubsection{馬太福音 2:2}
「那生下來做猶太人之王的在哪裡？我們在東方看見他的星，特來拜他。」

\subsubsection{路加福音 1:67-79}
67 他父親撒迦利亞被聖靈充滿了，就預言說： 68 「主以色列的神是應當稱頌的！因他眷顧他的百姓，為他們施行救贖， 69 在他僕人大衛家中，為我們興起了拯救的角， 70 正如主藉著從創世以來聖先知的口所說的話， 71 拯救我們脫離仇敵和一切恨我們之人的手； 72 向我們列祖施憐憫、記念他的聖約， 73 就是他對我們祖宗亞伯拉罕所起的誓， 74 叫我們既從仇敵手中被救出來， 75 就可以終身在他面前，坦然無懼地用聖潔、公義侍奉他。 76 孩子啊，你要稱為至高者的先知，因為你要行在主的前面，預備他的道路， 77 叫他的百姓因罪得赦就知道救恩。 78 因我們神憐憫的心腸，叫清晨的日光從高天臨到我們， 79 要照亮坐在黑暗中死蔭裡的人，把我們的腳引到平安的路上。」

\subsubsection{彼得後書 1:19}
我們並有先知更確的預言，如同燈照在暗處。你們在這預言上留意，直等到天發亮、晨星在你們心裡出現的時候，才是好的。

\subsubsection{啟示錄 2:27-29，22:16}
\begin{itemize} 
\itemsep-.25em
\item 他必用鐵杖轄管他們，將他們如同窯戶的瓦器打得粉碎，像我從我父領受的權柄一樣。我又要把晨星賜給他。聖靈向眾教會所說的話，凡有耳的，就應當聽！』
\item  22:16 「我耶穌差遣我的使者為眾教會將這些事向你們證明。我是大衛的根，又是他的後裔；我是明亮的晨星。」
\end{itemize} 


\subsection{成全}

\subsubsection{成全創世之工 創世記 2:2}
到第七日，神造物的工已經完畢，就在第七日歇了他一切的工，安息了。

\subsubsection{成全救赎之工 约19:30}
耶稣尝了那醋，就说，‘成了。’便低下头，将灵魂交付神了。

\subsubsection{神的奧祕成全、新造成全  啟示錄 10:1-7/21:6-8}
\begin{itemize}
\item 我又看見另有一位大力的天使從天降下，披著雲彩，頭上有虹，臉面像日頭，兩腳像火柱。他手裡拿著小書卷，是展開的。他右腳踏海，左腳踏地， 
大聲呼喊，好像獅子吼叫。呼喊完了，就有七雷發聲。七雷發聲之後，我正要寫出來，就聽見從天上有聲音說：七雷所說的，你要封上，不可寫出來。我所看見的那踏海踏地的天使向天舉起右手來，指著那創造天和天上之物，地和地上之物，海和海中之物，直活到永永遠遠的，起誓說：不再有時日了（或作：不再耽延了）。但在第七位天使吹號發聲的時候，神的奧祕就成全了，正如神所傳給他僕人眾先知的佳音。 
\item 他又對我說：「都成了！我是阿拉法，我是俄梅戛；我是初，我是終。我要將生命泉的水白白賜給那口渴的人喝。 7 得勝的，必承受這些為業：我要做他的神，他要做我的兒子。 8 唯有膽怯的、不信的、可憎的、殺人的、淫亂的、行邪術的、拜偶像的和一切說謊話的，他們的份就在燒著硫磺的火湖裡，這是第二次的死。」
\end{itemize}

\subsection{奥秘}
 


\subsubsection{天國的奧祕 太13:10-16/可4:10-12/路8:9-10}
\begin{itemize}
\item 門徒進前來，問耶穌說：「對眾人講話為什麼用比喻呢？」 11 耶穌回答說：「因為天國的奧祕只叫你們知道，不叫他們知道。 12 凡有的，還要加給他，叫他有餘；凡沒有的，連他所有的也要奪去。 13 所以我用比喻對他們講，是因他們看也看不見，聽也聽不見，也不明白。 14 在他們身上，正應了以賽亞的預言說：『你們聽是要聽見，卻不明白；看是要看見，卻不曉得。 15 因為這百姓油蒙了心，耳朵發沉，眼睛閉著，恐怕眼睛看見，耳朵聽見，心裡明白，回轉過來，我就醫治他們。』 16 但你們的眼睛是有福的，因為看見了；你們的耳朵也是有福的，因為聽見了。
\item 無人的時候，跟隨耶穌的人和十二個門徒問他這比喻的意思。 11 耶穌對他們說：「神國的奧祕只叫你們知道，若是對外人講，凡事就用比喻， 12 叫『他們看是看見，卻不曉得，聽是聽見，卻不明白；恐怕他們回轉過來，就得赦免』。」
\item 門徒問耶穌說：「這比喻是什麼意思呢？」 10 他說：「神國的奧祕只叫你們知道，至於別人，就用比喻，叫他們『看也看不見，聽也聽不明』。
\end{itemize}



\subsubsection{以色列人有幾分硬心的奧祕 羅11：25-32}
弟兄們，我不願意你們不知道這奧祕，恐怕你們自以為聰明，就是：以色列人有幾分是硬心的，等到外邦人的數目添滿了， 26 於是以色列全家都要得救。如經上所記：「必有一位救主從錫安出來，要消除雅各家的一切罪惡。」 27 又說：「我除去他們罪的時候，這就是我與他們所立的約。」 28 就著福音說，他們為你們的緣故是仇敵；就著揀選說，他們為列祖的緣故是蒙愛的。 29 因為神的恩賜和選召是沒有後悔的。 30 你們從前不順服神，如今因他們的不順服，你們倒蒙了憐恤。 31 這樣，他們也是不順服，叫他們因著施給你們的憐恤，現在也就蒙憐恤。 32 因為神將眾人都圈在不順服之中，特意要憐恤眾人。深哉，神豐富的智慧和知識！他的判斷何其難測，他的蹤跡何其難尋！ 34 「誰知道主的心，誰做過他的謀士呢？ 35 誰是先給了他，使他後來償還呢？」 36 因為萬有都是本於他，倚靠他，歸於他。願榮耀歸給他，直到永遠！阿們。


\subsubsection{福音和耶穌基督的奧祕 羅16：25-27}
唯有神能照我所傳的福音和所講的耶穌基督，並照永古隱藏不言的奧祕，堅固你們的心。 26 這奧祕如今顯明出來，而且按著永生神的命，藉眾先知的書指示萬國的民，使他們信服真道。 27 願榮耀因耶穌基督歸於獨一全智的神，直到永遠！阿們。


\subsubsection{各樣的奧祕 林前13:1-2/14:1-11}
\begin{itemize}
\item 我若能說萬人的方言，並天使的話語，卻沒有愛，我就成了鳴的鑼、響的鈸一般。 2 我若有先知講道之能，也明白\textcolor{red}{各樣的奧祕}、各樣的知識，而且有全備的信叫我能夠移山，卻沒有愛，我就算不得什麼。
\item 你們要追求愛，也要切慕屬靈的恩賜，其中更要羨慕的是做先知講道。 2 那說方言的原不是對人說，乃是對神說，因為沒有人聽出來，然而他在心靈裡卻是講說\textcolor{red}{各樣的奧祕}。 3 但做先知講道的是對人說，要造就、安慰、勸勉人。 4 說方言的是造就自己，做先知講道的乃是造就教會。 5 我願意你們都說方言，更願意你們做先知講道；因為說方言的若不翻出來，使教會被造就，那做先知講道的就比他強了。 6 弟兄們，我到你們那裡去，若只說方言，不用啟示或知識或預言或教訓給你們講解，我於你們有什麼益處呢？ 7 就是那有聲無氣的物，或簫或琴，若發出來的聲音沒有分別，怎能知道所吹所彈的是什麼呢？ 8 若吹無定的號聲，誰能預備打仗呢？ 9 你們也是如此，舌頭若不說容易明白的話，怎能知道所說的是什麼呢？這就是向空說話了。 10 世上的聲音或者甚多，卻沒有一樣是無意思的。 11 我若不明白那聲音的意思，這說話的人必以我為化外之人，我也以他為化外之人。
\end{itemize}


\subsubsection{奧祕的事 林前15:50-54}
弟兄們，我告訴你們說，血肉之體不能承受神的國，必朽壞的不能承受不朽壞的。 51 我如今把\textcolor{red}{一件奧祕的事}告訴你們：我們不是都要睡覺，乃是都要改變， 52 就在一霎時，眨眼之間，號筒末次吹響的時候。因號筒要響，死人要復活成為不朽壞的，我們也要改變。 53 這必朽壞的總要變成不朽壞的，這必死的總要變成不死的。 54 這必朽壞的既變成不朽壞的，這必死的既變成不死的，那時經上所記「死被得勝吞滅」的話就應驗了。


\subsubsection{神旨意的奧祕 弗1:3-14}
願頌讚歸於我們主耶穌基督的父神！他在基督裡曾賜給我們天上各樣屬靈的福氣， 4 就如神從創立世界以前在基督裡揀選了我們，使我們在他面前成為聖潔，無有瑕疵； 5 又因愛我們，就按著自己意旨所喜悅的，預定我們藉著耶穌基督得兒子的名分， 6 使他榮耀的恩典得著稱讚。這恩典是他在愛子裡所賜給我們的。 7 我們藉這愛子的血得蒙救贖，過犯得以赦免，乃是照他豐富的恩典。 8 這恩典是神用諸般智慧聰明，充充足足賞給我們的， 9 都是照他自己所預定的美意，叫我們知道他\textcolor{red}{旨意的奧祕}， 10 要照所安排的，在日期滿足的時候，使天上、地上一切所有的，都在基督裡面同歸於一。 11 我們也在他裡面得 a了基業，這原是那位隨己意行做萬事的，照著他旨意所預定的， 12 叫他的榮耀從我們這首先在基督裡有盼望的人，可以得著稱讚。 13 你們既聽見真理的道，就是那叫你們得救的福音，也信了基督，既然信他，就受了所應許的聖靈為印記。 14 這聖靈是我們得基業的憑據 b，直等到神之民 c被贖，使他的榮耀得著稱讚。


\subsubsection{福音的奧祕 弗3:1-9}
因此，我保羅——為你們外邦人做了基督耶穌被囚的，替你們祈禱。 2 諒必你們曾聽見神賜恩給我，將關切你們的職分託付我， 3 用啟示使我知道\textcolor{red}{福音的奧祕}，正如我以前略略寫過的。 4 你們念了，就能曉得我深知基督的奧祕。 5 這奧祕在以前的世代沒有叫人知道，像如今藉著聖靈啟示他的聖使徒和先知一樣。 6 這奧祕就是外邦人在基督耶穌裡，藉著福音，得以同為後嗣，同為一體，同蒙應許。 7 我做了這福音的執事，是照神的恩賜，這恩賜是照他運行的大能賜給我的。 8 我本來比眾聖徒中最小的還小，然而他還賜我這恩典，叫我把基督那測不透的豐富傳給外邦人， 9 又使眾人都明白，這歷代以來隱藏在創造萬物之神裡的奧祕是如何安排的， 10 為要藉著教會，使天上執政的、掌權的現在得知神百般的智慧。 11 這是照神從萬世以前在我們主基督耶穌裡所定的旨意。 12 我們因信耶穌，就在他裡面放膽無懼，篤信不疑地來到神面前。 13 所以，我求你們不要因我為你們所受的患難喪膽，這原是你們的榮耀。


\subsubsection{神的奧祕 林前2:1-16/啟10:5-7/11:15-19}
\begin{itemize}
\item 弟兄們，從前我到你們那裡去，並沒有用高言大智對你們宣傳神的奧祕。 2 因為我曾定了主意，在你們中間不知道別的，只知道耶穌基督並他釘十字架。 3 我在你們那裡，又軟弱，又懼怕，又甚戰兢。 4 我說的話、講的道，不是用智慧委婉的言語，乃是用聖靈和大能的明證， 5 叫你們的信不在乎人的智慧，只在乎神的大能。然而，在完全的人中，我們也講智慧，但不是這世上的智慧，也不是這世上有權有位將要敗亡之人的智慧。 7 我們講的，乃是從前所隱藏、神奧祕的智慧，就是神在萬世以前預定使我們得榮耀的。 8 這智慧，世上有權有位的人沒有一個知道的，他們若知道，就不把榮耀的主釘在十字架上了。 9 如經上所記：「神為愛他的人所預備的，是眼睛未曾看見，耳朵未曾聽見，人心也未曾想到的。」 10 只有神藉著聖靈向我們顯明了。因為聖靈參透萬事，就是神深奧的事也參透了。 11 除了在人裡頭的靈，誰知道人的事？像這樣，除了神的靈，也沒有人知道神的事。 12 我們所領受的，並不是世上的靈，乃是從神來的靈，叫我們能知道神開恩賜給我們的事。 13 並且我們講說這些事，不是用人智慧所指教的言語，乃是用聖靈所指教的言語，將屬靈的話解釋屬靈的事。然而，屬血氣的人不領會神聖靈的事，反倒以為愚拙；並且不能知道，因為這些事唯有屬靈的人才能看透。 15 屬靈的人能看透萬事，卻沒有一人能看透了他。 16 「誰曾知道主的心，去教導他呢？」但我們是有基督的心了。
\item 我所看見的那踏海踏地的天使向天舉起右手來， 6 指著那創造天和天上之物、地和地上之物、海和海中之物，直活到永永遠遠的，起誓說：「不再有時日了 a！ 7 但在第七位天使吹號發聲的時候，神的奧祕就成全了，正如神所傳給他僕人眾先知的佳音。」
\item 第七位天使吹號，天上就有大聲音說：「世上的國成了我主和主基督的國，他要做王，直到永永遠遠。」 16 在神面前、坐在自己位上的二十四位長老就面伏於地，敬拜神， 17 說：「昔在、今在的主神，全能者啊！我們感謝你，因你執掌大權做王了。 18 外邦發怒，你的憤怒也臨到了！審判死人的時候也到了！你的僕人眾先知和眾聖徒，凡敬畏你名的人，連大帶小得賞賜的時候也到了！你敗壞那些敗壞世界之人的時候也就到了！」當時，神天上的殿開了，在他殿中現出他的約櫃，隨後有閃電、聲音、雷轟、地震、大雹。
\end{itemize}

\subsubsection{神奧祕事的管家 林前4:1-5}
人應當以我們為基督的執事，為神奧祕事的管家。 2 所求於管家的，是要他有忠心。 3 我被你們論斷，或被別人論斷，我都以為極小的事，連我自己也不論斷自己。 4 我雖不覺得自己有錯，卻也不能因此得以稱義，但判斷我的乃是主。 5 所以，時候未到，什麼都不要論斷，只等主來，他要照出暗中的隱情，顯明人心的意念。那時，各人要從神那裡得著稱讚。


 \subsubsection{七星和七個金燈臺的奧祕 啟示錄 1:20}
論到你所看見在我右手中的七星和七個金燈臺的奧祕，那七星就是七個教會的使者，七燈臺就是七個教會。
 


  which has been hidden

Ephesians 5:32 N-NNS
GRK: τὸ μυστήριον τοῦτο μέγα
NAS: This mystery is great;
KJV: is a great mystery: but I
INT: the mystery This great

Ephesians 6:19 N-ANS
GRK: γνωρίσαι τὸ μυστήριον τοῦ εὐαγγελίου
NAS: with boldness the mystery of the gospel,
KJV: to make known the mystery of the gospel,
INT: to make known the mystery of the gospel

Colossians 1:26 N-ANS
GRK: τὸ μυστήριον τὸ ἀποκεκρυμμένον
NAS: [that is], the mystery which has been hidden
KJV: [Even] the mystery which hath been hid
INT: the mystery which has been hidden

Colossians 1:27 N-GNS
GRK: δόξης τοῦ μυστηρίου τούτου ἐν
NAS: of this mystery among
KJV: of this mystery among
INT: glory of the mystery this among

Colossians 2:2 N-GNS
GRK: ἐπίγνωσιν τοῦ μυστηρίου τοῦ θεοῦ
NAS: of God's mystery, [that is], Christ
KJV: the acknowledgement of the mystery of God,
INT: [the] knowledge of the mystery of God

Colossians 4:3 N-ANS
GRK: λαλῆσαι τὸ μυστήριον τοῦ χριστοῦ
NAS: so that we may speak forth the mystery of Christ,
KJV: to speak the mystery of Christ,
INT: to speak the mystery of Christ

2 Thessalonians 2:7 N-NNS
GRK: τὸ γὰρ μυστήριον ἤδη ἐνεργεῖται
NAS: For the mystery of lawlessness
KJV: For the mystery of iniquity doth
INT: the indeed mystery already is working

1 Timothy 3:9 N-ANS
GRK: ἔχοντας τὸ μυστήριον τῆς πίστεως
NAS: [but] holding to the mystery of the faith
KJV: Holding the mystery of the faith in
INT: holding the mystery of the faith

1 Timothy 3:16 N-NNS
GRK: τῆς εὐσεβείας μυστήριον Ὃς ἐφανερώθη
NAS: great is the mystery of godliness:
KJV: great is the mystery of godliness: God
INT: of godliness mystery which was revealed

Revelation 1:20 N-NNS 



 
 

 

 
 
 
 
 
 

 
    
     
 

 
 


 

     
     

 %\subsubsection{羅馬書 2:1-11}    
%你這論斷人的，無論你是誰，也無可推諉。你在什麼事上論斷人，就在什麼事上定自己的罪，因你這論斷人的，自己所行卻和別人一樣。 2 我們知道這樣行的人，神必照真理審判他。 3 你這人哪，你論斷行這樣事的人，自己所行的卻和別人一樣，你以為能逃脫神的審判嗎？ 4 還是你藐視他豐富的恩慈、寬容、忍耐，不曉得他的恩慈是領你悔改呢？ 5 你竟任著你剛硬不悔改的心，為自己積蓄憤怒，以致神震怒，顯他公義審判的日子來到。 6 他必照各人的行為報應各人。 7 凡恆心行善，尋求榮耀、尊貴和不能朽壞之福的，就以永生報應他們； 8 唯有結黨、不順從真理反順從不義的，就以憤怒、惱恨報應他們。 9 將\underline{患難$\theta$$\lambda$$\tilde{\iota}$$\psi$$\iota$$\zeta$θλῖψις}、困苦加給一切作惡的人，先是猶太人，後是希臘人； 10 卻將榮耀、尊貴、平安加給一切行善的人，先是猶太人，後是希臘人。 11 因為神不偏待人。  
   


   
 


 

    
    
    
  


 
   
  
      
 
  
                

 

 


\subsection{金燈臺和二橄欖樹指二受膏者}

\subsubsection{啟示錄11：3-13}
 3 我要使我那兩個見證人穿著毛衣，傳道一千二百六十天。」 4 他們就是那兩棵橄欖樹，兩個燈臺，立在世界之主面前的。 5 若有人想要害他們，就有火從他們口中出來，燒滅仇敵。凡想要害他們的都必這樣被殺。 6 這二人有權柄在他們傳道的日子叫天閉塞不下雨，又有權柄叫水變為血，並且能隨時隨意用各樣的災殃攻擊世界。 7 他們作完見證的時候，那從無底坑裡上來的獸必與他們交戰，並且得勝，把他們殺了。8 他們的屍首就倒在大城裡的街上，這城按著靈意叫所多瑪，又叫埃及，就是他們的主釘十字架之處。 9 從各民、各族、各方、各國中，有人觀看他們的屍首三天半，又不許把屍首放在墳墓裡。 10 住在地上的人就為他們歡喜快樂，互相饋送禮物，因這兩位先知曾叫住在地上的人受痛苦。過了這三天半，有生氣從神那裡進入他們裡面，他們就站起來，看見他們的人甚是害怕。 12 兩位先知聽見有大聲音從天上來，對他們說：「上到這裡來！」他們就駕著雲上了天，他們的仇敵也看見了。 13 正在那時候，地大震動，城就倒塌了十分之一，因地震而死的有七千人，其餘的都恐懼，歸榮耀給天上的神。

\subsubsection{撒迦利亞書 4：1-14}
那與我說話的天使又來叫醒我，好像人睡覺被喚醒一樣。 2 他問我說：「你看見了什麼？」我說：「我看見了一個純金的燈臺，頂上有燈盞，燈臺上有七盞燈，每盞有七個管子。 3 旁邊有兩棵橄欖樹，一棵在燈盞的右邊，一棵在燈盞的左邊。」 4 我問與我說話的天使說：「主啊，這是什麼意思？」 5 與我說話的天使回答我說：「你不知道這是什麼意思嗎？」我說：「主啊，我不知道。」 6 他對我說：「這是耶和華指示所羅巴伯的。萬軍之耶和華說：不是倚靠勢力，不是倚靠才能，乃是倚靠我的靈方能成事。 7 大山哪，你算什麼呢？在所羅巴伯面前，你必成為平地。他必搬出一塊石頭，安在殿頂上，人且大聲歡呼說：『願恩惠，恩惠歸於這殿 a！』」 8 耶和華的話又臨到我說： 9 「所羅巴伯的手立了這殿的根基，他的手也必完成這工，你就知道萬軍之耶和華差遣我到你們這裡來了。 10 誰藐視這日的事為小呢？這七眼乃是耶和華的眼睛，遍察全地，見所羅巴伯手拿線砣就歡喜。」
我又問天使說：「這燈臺左右的兩棵橄欖樹是什麼意思？」 12 我二次問他說：「這兩根橄欖枝在兩個流出金色油的金嘴旁邊是什麼意思？」 13 他對我說：「你不知道這是什麼意思嗎？」我說：「主啊，我不知道。」 14 他說：「這是兩個受膏者，站在普天下主的旁邊。」

\subsection{婦人}
\subsubsection{身披日頭，腳踏月亮，頭戴十二星冠冕的婦人}
天上現出大異象來：有一個婦人身披日頭，腳踏月亮，頭戴十二星的冠冕， 2 她懷了孕，在生產的艱難中疼痛呼叫。
\subsubsection{婦人的孩子被提到神寶座}
天上又現出異象來：有一條大紅龍，七頭十角，七頭上戴著七個冠冕， 4 牠的尾巴拖拉著天上星辰的三分之一，摔在地上。龍就站在那將要生產的婦人面前，等她生產之後，要吞吃她的孩子。 5 婦人生了一個男孩子，是將來要用鐵杖轄管 a萬國的；她的孩子被提到神寶座那裡去了。 6 婦人就逃到曠野，在那裡有神給她預備的地方，使她被養活一千二百六十天。

\subsubsection{撫養婦人一載二載半載}

13 龍見自己被摔在地上，就逼迫那生男孩子的婦人。 14 於是有大鷹的兩個翅膀賜給婦人，叫她能飛到曠野，到自己的地方躲避那蛇，她在那裡被養活一載，二載，半載。 15 蛇就在婦人身後，從口中吐出水來，像河一樣，要將婦人沖去。 16 地卻幫助婦人，開口吞了從龍口吐出來的水 b。 17 龍向婦人發怒，去與她其餘的兒女爭戰，這兒女就是那守神誡命、為耶穌作見證的。那時龍就站在海邊的沙上。


\subsection{龍}
\subsubsection{大紅龍要吞婦人的孩子}

天上又現出異象來：有一條大紅龍，七頭十角，七頭上戴著七個冠冕， 4 牠的尾巴拖拉著天上星辰的三分之一，摔在地上。龍就站在那將要生產的婦人面前，等她生產之後，要吞吃她的孩子。

\subsubsection{天上有爭戰龍就被摔於地}

7 在天上就有了爭戰，米迦勒同他的使者與龍爭戰。龍也同牠的使者去爭戰， 8 並沒有得勝，天上再沒有牠們的地方。 9 大龍就是那古蛇，名叫魔鬼，又叫撒旦，是迷惑普天下的。牠被摔在地上，牠的使者也一同被摔下去。 10 我聽見在天上有大聲音說：「我神的救恩、能力、國度並他基督的權柄，現在都來到了！因為那在我們神面前晝夜控告我們弟兄的，已經被摔下去了。 11 弟兄勝過牠，是因羔羊的血和自己所見證的道，他們雖至於死也不愛惜性命。 12 所以諸天和住在其中的，你們都快樂吧！只是地與海有禍了！因為魔鬼知道自己的時候不多，就氣憤憤地下到你們那裡去了。」



\subsection{兽}
\subsubsection{海中的兽}

我又看見一個獸從海中上來，有十角七頭，在十角上戴著十個冠冕，七頭上有褻瀆的名號。 2 我所看見的獸形狀像豹，腳像熊的腳，口像獅子的口。那龍將自己的能力、座位和大權柄都給了牠。 3 我看見獸的七頭中，有一個似乎受了死傷，那死傷卻醫好了。全地的人都稀奇，跟從那獸， 4 又拜那龍，因為牠將自己的權柄給了獸，也拜獸，說：「誰能比這獸，誰能與牠交戰呢？」 5 又賜給牠說誇大褻瀆話的口，又有權柄賜給牠，可以任意而行四十二個月。 6 獸就開口向神說褻瀆的話，褻瀆神的名並他的帳幕，以及那些住在天上的。又任憑牠與聖徒爭戰，並且得勝；也把權柄賜給牠，制伏各族、各民、各方、各國。 8 凡住在地上，名字從創世以來沒有記在被殺之羔羊生命冊上的人，都要拜牠。 9 凡有耳的，就應當聽！ 10 擄掠人的必被擄掠，用刀殺人的必被刀殺。聖徒的忍耐和信心就是在此。

\subsubsection{女人所騎之獸的奧祕}
我就看見一個女人騎在朱紅色的獸上，那獸有七頭十角，遍體有褻瀆的名號。你所看見的獸，先前有，如今沒有，將要從無底坑裡上來，又要歸於沉淪。凡住在地上，名字從創世以來沒有記在生命冊上的，見先前有、如今沒有、以後再有的獸，就必稀奇。 9 智慧的心在此可以思想。那七頭就是女人所坐的七座山， 10 又是七位王：五位已經傾倒了，一位還在，一位還沒有來到；他來的時候，必須暫時存留。 11 那先前有、如今沒有的獸，就是第八位，他也和那七位同列，並且歸於沉淪。 12 你所看見的那十角就是十王，他們還沒有得國，但他們一時之間要和獸同得權柄，與王一樣。 13 他們同心合意將自己的能力、權柄給那獸。「他們與羔羊爭戰，羔羊必勝過他們，因為羔羊是萬主之主、萬王之王。同著羔羊的，就是蒙召、被選、有忠心的，也必得勝。」

天使又對我說：「你所看見那淫婦坐的眾水，就是多民、多人、多國、多方。 16 你所看見的那十角與獸必恨這淫婦，使她冷落赤身，又要吃她的肉，用火將她燒盡。



%\subsection{地中的兽 }
\subsubsection{地中的兽}
我又看見另有一個獸從地中上來，有兩角如同羊羔，說話好像龍。 12 牠在頭一個獸面前，施行頭一個獸所有的權柄，並且叫地和住在地上的人拜那死傷醫好的頭一個獸。 13 又行大奇事，甚至在人面前叫火從天降在地上。 14 牠因賜給牠權柄在獸面前能行奇事，就迷惑住在地上的人，說要給那受刀傷還活著的獸做個像。 15 又有權柄賜給牠，叫獸像有生氣，並且能說話，又叫所有不拜獸像的人都被殺害。牠又叫眾人，無論大小，貧富，自主的為奴的，都在右手上或是在額上受一個印記。 17 除了那受印記，有了獸名或有獸名數目的，都不得做買賣。 18 在這裡有智慧：凡有聰明的，可以算計獸的數目，因為這是人的數目，它的數目是六百六十六。
\subsubsection{敵基督 2:18-25}
\begin{itemize}%[noitemsep]
\itemsep-.5em
\item  \underline{末時}那敵基督的要來
\begin{itemize}%[noitemsep]
\itemsep-.5em 
\item  現在已經有好些敵基督的出來了
\item  從此我們就知道如今是末時了
\end{itemize} 
\item  他們從我們中間出去，卻不是屬我們的;他們出去，顯明都不是屬我們的
\item  若是屬我們的，就必仍舊與我們同在
\item  不認父與子
\begin{itemize}%[noitemsep]
\itemsep-.5em 
\item  不認耶穌為基督的是說謊話的
\item  不認父與子的，這就是敵基督的
\item  凡不認子的，就沒有父；認子的，連父也有了
\item  要將那從起初所聽見的，常存在心裡。就必住在子裡面，也必住在父裡面
\end{itemize}
\end{itemize}


\subsection{大淫妇}
\subsubsection{大淫婦的刑罰 }
拿著七碗的七位天使中，有一位前來對我說：「你到這裡來，我將坐在眾水上的大淫婦所要受的刑罰指給你看。 2 地上的君王與她行淫，住在地上的人喝醉了她淫亂的酒。」 3 我被聖靈感動，天使帶我到曠野去。我就看見一個女人騎在朱紅色的獸上，那獸有七頭十角，遍體有褻瀆的名號。 4 那女人穿著紫色和朱紅色的衣服，用金子、寶石、珍珠為裝飾，手拿金杯，杯中盛滿了可憎之物，就是她淫亂的汙穢。 5 在她額上有名寫著說：「奧祕哉，大巴比倫，做世上的淫婦和一切可憎之物的母！」 6 我又看見那女人喝醉了聖徒的血和為耶穌作見證之人的血。我看見她，就大大地稀奇。
\subsubsection{大淫婦騎獸的奧祕}
天使又對我說：「你所看見那淫婦坐的眾水，就是多民、多人、多國、多方。 16 你所看見的那十角與獸必恨這淫婦，使她冷落赤身，又要吃她的肉，用火將她燒盡。 17 因為神使諸王同心合意，遵行他的旨意，把自己的國給那獸，直等到神的話都應驗了。 18 你所看見的那女人就是管轄地上眾王的大城。」


\subsubsection{巴比倫傾倒}

1 此後，我看見另有一位有大權柄的天使從天降下，地就因他的榮耀發光。 2 他大聲喊著說：「巴比倫大城傾倒了！傾倒了！成了鬼魔的住處和各樣汙穢之靈的巢穴 a，並各樣汙穢可憎之雀鳥的巢穴。 3 因為列國都被她邪淫大怒的酒傾倒了，地上的君王與她行淫，地上的客商因她奢華太過就發了財。」

\subsubsection{先有榮耀奢華後受悲哀痛苦}
我又聽見從天上有聲音說：「我的民哪，你們要從那城出來，免得與她一同有罪，受她所受的災殃。 5 因她的罪惡滔天，她的不義，神已經想起來了。 6 她怎樣待人，也要怎樣待她，按她所行的加倍地報應她，用她調酒的杯加倍地調給她喝。 7 她怎樣榮耀自己，怎樣奢華，也當叫她照樣痛苦悲哀；因她心裡說：『我坐了皇后的位，並不是寡婦，決不至於悲哀。』 8 所以在一天之內，她的災殃要一齊來到，就是死亡、悲哀、饑荒。她又要被火燒盡了，因為審判她的主神大有能力。 9 地上的君王，素來與她行淫、一同奢華的，看見燒她的煙，就必為她哭泣哀號。 10 因怕她的痛苦，就遠遠地站著說：『哀哉！哀哉！巴比倫大城，堅固的城啊，一時之間你的刑罰就來到了！』 11 地上的客商也都為她哭泣悲哀，因為沒有人再買他們的貨物了。 12 這貨物就是金、銀、寶石、珍珠、細麻布、紫色料、綢子、朱紅色料，各樣香木，各樣象牙的器皿，各樣極寶貴的木頭和銅、鐵、漢白玉的器皿， 13 並肉桂、豆蔻、香料、香膏、乳香、酒、油、細麵、麥子、牛、羊、車、馬和奴僕、人口。

\subsubsection{因大城受報應聖徒都歡喜}
「巴比倫哪，你所貪愛的果子離開了你，你一切的珍饈美味和華美的物件也從你中間毀滅，決不能再見了！ 15 販賣這些貨物，藉著她發了財的客商，因怕她的痛苦，就遠遠地站著哭泣悲哀， 16 說：『哀哉，哀哉，這大城啊！素常穿著細麻、紫色、朱紅色的衣服，又用金子、寶石和珍珠為裝飾， 17 一時之間，這麼大的富厚就歸於無有了！』凡船主和坐船往各處去的，並眾水手，連所有靠海為業的，都遠遠地站著， 18 看見燒她的煙就喊著說：『有何城能比這大城呢？』 19 他們又把塵土撒在頭上，哭泣悲哀，喊著說：『哀哉，哀哉，這大城啊！凡有船在海中的都因她的珍寶成了富足，她在一時之間就成了荒場！』 20 天哪，眾聖徒、眾使徒、眾先知啊！你們都要因她歡喜，因為神已經在她身上申了你們的冤。」

21 有一位大力的天使舉起一塊石頭，好像大磨石，扔在海裡說：「巴比倫大城也必這樣猛力地被扔下去，決不能再見了。 22 彈琴、作樂、吹笛、吹號的聲音在你中間決不能再聽見。各行手藝人在你中間決不能再遇見。推磨的聲音在你中間決不能再聽見。 23 燈光在你中間決不能再照耀。新郎和新婦的聲音在你中間決不能再聽見。你的客商原來是地上的尊貴人，萬國也被你的邪術迷惑了。 24 先知和聖徒並地上一切被殺之人的血，都在這城裡看見了。」

\subsubsection{神申討大淫婦流人血之罪}

1 此後，我聽見好像群眾在天上大聲說：「哈利路亞 a！救恩、榮耀、權能都屬乎我們的神！ 2 他的判斷是真實、公義的，因他判斷了那用淫行敗壞世界的大淫婦，並且向淫婦討流僕人血的罪，給他們申冤。」 3 又說：「哈利路亞！燒淫婦的煙往上冒，直到永永遠遠！」 4 那二十四位長老與四活物就俯伏敬拜坐寶座的神，說：「阿們！哈利路亞！」 5 有聲音從寶座出來，說：「神的眾僕人哪，凡敬畏他的，無論大小，都要讚美我們的神！」 6 我聽見好像群眾的聲音，眾水的聲音，大雷的聲音，說：「哈利路亞！因為主我們的神，全能者做王了。




























 


















